\chapter{Universe levels and utilities}
\label{chap:levels}

This chapter covers four verification files that sit beside the main typing development.
\srcloc{Lean4Lean/Verify/Level.lean}{1} (3880 lines) verifies lean4lean's own universe-level
canonical form; \srcloc{Lean4Lean/Verify/NormLt.lean}{1} (364 lines) verifies the order core Lean
uses to sort the arguments of a \texttt{max}; \srcloc{Lean4Lean/Verify/QSort.lean}{1} (392 lines)
is a vendored specification of \texttt{Array.qsort} that the previous file needs; and
\srcloc{Lean4Lean/Verify/EquivManager.lean}{1} (331 lines) proves the kernel's memo cache for
structural equality sound. Every line of all four files is attributed to master: the
\texttt{iota} and \texttt{trproj} branches contribute nothing here (0 lines in each of the four
files), so this chapter has no contributed nodes and exists to cover what the contributions rest
on. There is no active \texttt{sorry} in the group; the one occurrence of the token sits inside a
dead comment block at \srcloc{Lean4Lean/Verify/Level.lean}{27}.

\texttt{Lean4Lean/Level.lean} implements a decision procedure for level equality following Yoan
Géran, \emph{A Canonical Form for Universe Levels in Impredicative Type Theory}, and
\srcloc{Lean4Lean/Verify/Level.lean}{166} proves it correct. A level is represented as a
\texttt{NormLevel}, a \texttt{Std.TreeMap} from \emph{condition sets} (lists of parameter names,
the guards of an \texttt{imax} chain) to nodes; the node at a key $p$ bundles a constant sublevel
$C(p,c)$ and a list of variable sublevels $V(p,x+k)$, in the notation of the paper
(\srcloc{Lean4Lean/Level.lean}{41}). The verification runs in six phases: preliminaries about
core \texttt{Level} (lines 1--165), accumulation and semantics (166--660), subsumption
(660--1290), reconstruction of a \texttt{Lean.Level} from a normal form (1290--2250),
completeness and canonicity (2250--2940, resuming at 3397--3792 after the flat-path material),
and a flat fast path (2940--3396 and 3793--3880). Throughout, $ls$ is
the list of universe parameter names in scope and $\rho$ a list of naturals assigning them
values; \texttt{NormLevel.eval} $ls$ $\rho$ denotes a normal form and \texttt{VLevel.eval}
(\ref{def:vlevel-eval}) denotes a model level, so every soundness statement of this chapter is an
equation or an inequality between natural numbers, for all $\rho$.

\thesisref{axioms.tex, the algorithmic level inequality $\ell \le \ell'$ and the equivalence
$\ell \equiv \ell'$ it defines (rules only, no algorithm); soundness.tex, the valuation semantics
of levels} The thesis specifies level equality by a syntactic derivation system and gives the
semantics informally; it says nothing about deciding either. lean4lean instead \emph{defines}
level inequality semantically (\ref{def:vlevel-eval}) and proves a decision procedure sound and
complete against it, so the canonical form, its canonicity and the completeness results below are
a genuine extension beyond the thesis rather than a formalisation of it. Nothing in
\srcloc{Lean4Lean/Verify/NormLt.lean}{1}, \srcloc{Lean4Lean/Verify/QSort.lean}{1} or
\srcloc{Lean4Lean/Verify/EquivManager.lean}{1} corresponds to the thesis at all: those are
implementation-fidelity obligations arising from Lean's actual kernel (sorting inside
\texttt{Level.normalize}, and the \texttt{isDefEq} memo cache).

The trust boundary is worth stating up front. \texttt{NormLevel.le}, \texttt{NormLevel.addable}
and the \texttt{BEq} instance on \texttt{NormLevel} are all defined through
\texttt{Std.TreeMap.all}/\texttt{any}, so every main theorem about normal forms factors through
the two axioms \texttt{all\_eq\_all\_toList} and \texttt{any\_eq\_any\_toList}
(\ref{family:treemap-axioms}, \srcloc{Lean4Lean/Verify/Axioms.lean}{10}). The results that reach
core Lean's own \texttt{normalize} additionally inherit the axioms of
\srcloc{Lean4Lean/Verify/LevelStd.lean}{1} (\ref{family:level-axioms}), and the bit-packing
lemmas inherit \texttt{Lean.Level.mkData\_eq} (\srcloc{Lean4Lean/Verify/Axioms.lean}{269}). The
equivalence manager rests on \texttt{ptrEqExpr\_eq} (\ref{axiom:kernel-ptr-eq}) and, for the
\texttt{proj} case, on the still open \texttt{TrProj.uniq} (\ref{thm:trproj-uniq}), which is the
only route by which \texttt{sorryAx} enters this group.

Overall assessment: this is upstream work of uniformly high quality, and
\srcloc{Lean4Lean/Verify/Level.lean}{1} is the strongest single artefact in it. The statements
are the right ones (semantic, not syntactic), the section docstrings explain design choices
rather than restating code, and \texttt{NormLevel.separation} (\ref{thm:lvl-separation}) is a
real mathematical theorem proved without automation black boxes. Its weaknesses are the ordinary
ones of a large research file: a handful of dead lemmas that drag in an otherwise unused axiom,
a stale test-file docstring, five or six invariants threaded by hand through every end-to-end
statement, and a substantial verified component (the whole reconstruction layer supporting
\texttt{normalize'}) that no kernel code path yet calls. The two smaller level files are clean
but lean on \texttt{grind} in places that a reader cannot then audit, and
\srcloc{Lean4Lean/Verify/QSort.lean}{1} is vendored third-party code that should be replaced by a
standard-library specification when one lands.

\section{Preliminaries on core \texttt{Level}}

The first 165 lines of \srcloc{Lean4Lean/Verify/Level.lean}{1} are unrelated to the canonical
form: they are facts about core Lean's \texttt{Level} API that other verification files need.

\begin{definition}[Structural offset of a level]
\label{def:lvl-getoffset}
\lean{Lean.Level.getOffset', Lean.Level.getOffset_eq}
\leanok
\srcloc{Lean4Lean/Verify/Level.lean}{29}
\texttt{getOffset'} counts the leading \texttt{succ} constructors of a level by structural
recursion, and \texttt{getOffset\_eq} shows that it agrees with core Lean's accumulator-based
\texttt{Level.getOffset}. The proof is the generalised identity
\texttt{getOffsetAux u i = getOffset' u + i}, by induction on \texttt{u}.
\end{definition}

\begin{lemma}[Bit-packing of level metadata]
\label{fam:lvl-mkdata}
\lean{Lean.Level.mkData_depth, Lean.Level.mkData_hasParam, Lean.Level.mkData_hasMVar}
\leanok
\srcloc{Lean4Lean/Verify/Level.lean}{41}
For a depth $d < 2^{24}$, the packed \texttt{Level.Data} word \texttt{mkData h d hmv hp} reads
back the depth $d$, the \texttt{hasParam} flag \texttt{hp} and the \texttt{hasMVar} flag
\texttt{hmv} unchanged.
\end{lemma}
\begin{proof}
\leanok
Rewrite with the axiom \texttt{Lean.Level.mkData\_eq} (\srcloc{Lean4Lean/Verify/Axioms.lean}{269})
into the model \texttt{mkData'}, discharge the depth guard, and reduce the goal to a bit-vector
identity on \texttt{UInt64} closed by \texttt{bv\_decide}. Elaborating these three statements
requires \texttt{set\_option allowUnsafeReducibility true} together with
\texttt{attribute [local reducible] Data} (lines 38--39).
\axfoot{Lean.Level.mkData\_eq; three \texttt{bv\_decide} LRAT certificate axioms (\_native.bv\_decide.ax\_1\_8/9)}
\end{proof}

\begin{theorem}[Parameter-free, metavariable-free levels always translate]
\label{thm:lvl-oflevel-of-not-hasparam}
\lean{Lean.Level.ofLevel_of_not_hasParam}
\leanok
\srcloc{Lean4Lean/Verify/Level.lean}{88}
\uses{def:vlevel-oflevel}
If \texttt{l.hasParam'} and \texttt{l.hasMVar'} are both \texttt{false}, then
\texttt{VLevel.ofLevel Us l} succeeds for every list of universe names \texttt{Us}.
\end{theorem}
\begin{proof}
\leanok
Induction on \texttt{l}, then \texttt{simp\_all} with the definitions of \texttt{hasParam'},
\texttt{hasMVar'} and \texttt{VLevel.ofLevel}.
\end{proof}

\begin{theorem}[No undefined parameter implies translatability]
\label{thm:lvl-getundefparam-none}
\lean{Lean.Level.getUndefParam_none, Lean.Level.getUndefParam.F}
\leanok
\srcloc{Lean4Lean/Verify/Level.lean}{101}
\uses{def:kernel-level-getundefparam, def:vlevel-oflevel}
If \texttt{l} has no metavariables and \texttt{l.getUndefParam Us = none}, that is, no parameter
outside \texttt{Us} occurs in \texttt{l}, then \texttt{VLevel.ofLevel Us l} succeeds.
\texttt{getUndefParam.F} is the model of the monadic step that core's \texttt{Level.forEach}
runs.
\end{theorem}
\begin{proof}
\uses{thm:lvl-oflevel-of-not-hasparam}
\leanok
Strengthen the claim to a statement about the state-monad run of \texttt{Level.forEach}: for
every starting state \texttt{s}, a final state of \texttt{none} forces \texttt{s = none} and
translatability. Then induct on the level, handling the early-exit behaviour of the accumulator
explicitly; the \texttt{param} case appeals to
\texttt{ofLevel\_of\_not\_hasParam}. This is the form consumed by
\srcloc{Lean4Lean/Verify/TypeChecker/InferType.lean}{31}.
\end{proof}

\begin{definition}[Model of level-parameter substitution]
\label{def:lvl-substparams}
\lean{Lean.Level.substParams', Lean.Level.substParams_eq, Lean.Level.substParams_eq_self, Lean.Level.substParams_id}
\leanok
\srcloc{Lean4Lean/Verify/Level.lean}{138}
\texttt{substParams' s red u} is a structurally recursive model of core Lean's
\texttt{Level.substParams}: it replaces parameters using \texttt{s} and, when the \texttt{red}
flag says that a parameter actually occurs below, rebuilds \texttt{max}/\texttt{imax} with the
smart constructors. \texttt{substParams\_eq} proves the model agrees with core (opening the
private \texttt{substParams.go} and inducting, splitting on the \texttt{hasParam} guard);
\texttt{substParams\_eq\_self} that substitution is the identity on parameter-free levels; and
\texttt{substParams\_id} that substituting each parameter for itself without reduction is the
identity. Consumed by \srcloc{Lean4Lean/Verify/Expr.lean}{1128} and
\srcloc{Lean4Lean/Verify/Typing/Lemmas.lean}{1623}.
\end{definition}

\section{The canonical form and its semantics}

\begin{definition}[Semantics of the normal form]
\label{def:lvl-normlevel-eval}
\lean{Lean.Level.Normalize.evalParam, Lean.Level.Normalize.VarNode.eval, Lean.Level.Normalize.Node.eval, Lean.Level.Normalize.allNZ, Lean.Level.Normalize.evalPath, Lean.Level.Normalize.NormLevel.eval}
\leanok
\srcloc{Lean4Lean/Verify/Level.lean}{234}
\uses{def:vlevel-eval, def:kernel-level-normalize}
The denotation of a normal form under a valuation. \texttt{evalParam ls} $\rho$ \texttt{x} looks
the name \texttt{x} up in \texttt{ls} and reads its value off $\rho$ (returning $0$ off the
list); a \texttt{VarNode} $x+k$ evaluates to \texttt{evalParam}$(x) + k$; a \texttt{Node}
evaluates to the maximum of its constant and its variable nodes; \texttt{allNZ path} says every
name in a condition set has a nonzero value; \texttt{evalPath path n} is \texttt{n} when
\texttt{allNZ path} holds and $0$ otherwise, which is the semantics of an \texttt{imax} guard
chain; and a whole \texttt{NormLevel} evaluates to the maximum over its entries of
\texttt{evalPath p (Node.eval n)}. This is the interpretation of the paper's sublevel functions
$C(S,k)$ and $V(S,v,k)$.
\thesisref{soundness.tex, the valuation semantics of levels; \texttt{evalPath} is exactly the
\texttt{imax} clause generalised from one guard to a set of guards}
\end{definition}

\begin{lemma}[Boundedness characterisations of the semantics]
\label{fam:lvl-eval-le}
\lean{Lean.Level.Normalize.Node.eval_le, Lean.Level.Normalize.NormLevel.eval_le, Lean.Level.Normalize.evalPath_cons, Lean.Level.Normalize.evalPath_max, Lean.Level.Normalize.evalPath_mono, Lean.Level.Normalize.evalPath_le, Lean.Level.Normalize.allNZ_cons, Lean.Level.Normalize.allNZ_mono}
\leanok
\srcloc{Lean4Lean/Verify/Level.lean}{247}
The family of $x \le n$ characterisations that drives almost every later proof: a node is
bounded by $n$ exactly when its constant and each of its variable nodes are; a normal form is
bounded by $n$ exactly when every entry's \texttt{evalPath} contribution is; \texttt{evalPath}
commutes with \texttt{max}, is monotone in its payload, and satisfies
\texttt{evalPath path n} $\le m$ iff \texttt{allNZ path} implies $n \le m$; and \texttt{allNZ} is
antitone in the condition set.
\end{lemma}
\begin{proof}
\uses{def:lvl-normlevel-eval}
\leanok
The two \texttt{eval\_le} lemmas unfold the \texttt{foldl} into a \texttt{foldr} over the
reversed list and induct, using \texttt{Nat.max\_le}; the \texttt{evalPath} lemmas unfold the
definition and split on the \texttt{allNZ} guard. The whole development downstream is written in
this $\le m$ style, with \texttt{ext\_le}/\texttt{le\_ext\_le} (lines 345--348) converting back to
equalities, which keeps the statements first-order at the price of indirection.
\end{proof}

\begin{definition}[Inserting one element into a condition set]
\label{def:lvl-extend}
\lean{Lean.Level.Normalize.Extend1, Lean.Level.Normalize.Extend?, Lean.Level.Normalize.Extend1.orderedInsert, Lean.Level.Normalize.Extend?.orderedInsert}
\leanok
\srcloc{Lean4Lean/Verify/Level.lean}{186}
\texttt{Extend1 l v l'} says that \texttt{l'} is \texttt{l} with \texttt{v} inserted at some
position; \texttt{Extend? l v l'} additionally allows \texttt{l' = l} when \texttt{v} was already
present. The two \texttt{orderedInsert} lemmas connect them to the algorithm's
\texttt{orderedInsert}, which returns \texttt{none} exactly when the element is already there.
\end{definition}

\begin{definition}[Strictly sorted condition sets]
\label{def:lvl-sorted}
\lean{Lean.Level.Normalize.Sorted, Lean.Level.Normalize.Sorted.orderedInsert, Lean.Level.Normalize.Sorted.nodup, Lean.Level.Normalize.Sorted.erase, Lean.Level.Normalize.Sorted.head, Lean.Level.Normalize.Sorted.of_cons}
\leanok
\srcloc{Lean4Lean/Verify/Level.lean}{352}
\uses{def:lvl-extend}
\texttt{Sorted l} says the list of names \texttt{l} is pairwise strictly increasing in
\texttt{Name.cmp}; in particular it is duplicate-free. This is the invariant that makes the
linear merge \texttt{subset} a decision procedure for set inclusion, and it is preserved by
\texttt{orderedInsert} and by \texttt{List.erase}.
\end{definition}

\begin{definition}[Variable lists are sorted by name]
\label{def:lvl-varssorted}
\lean{Lean.Level.Normalize.VarsSorted, Lean.Level.Normalize.NormLevel.SortedVars, Lean.Level.Normalize.VarNode.addVar_sorted, Lean.Level.Normalize.normalizeAux_sortedVars, Lean.Level.Normalize.NormLevel.subsumption_sortedVars, Lean.Level.Normalize.VarsSorted.eq_of_var_eq, Lean.Level.Normalize.VarsSorted.eq_of_mem_iff}
\leanok
\srcloc{Lean4Lean/Verify/Level.lean}{2263}
\uses{def:lvl-sorted}
\texttt{VarsSorted l} says a \texttt{VarNode} list is strictly increasing in the variable name,
so each name occurs at most once, and \texttt{SortedVars s} lifts this to every entry of a map.
The invariant is established by \texttt{VarNode.addVar}, preserved by \texttt{normalizeAux} and
by subsumption, and makes two variable lists with the same members equal. It is needed twice
downstream: for the exactness of \texttt{subsumeVars} (a surviving variable really has no
dominator) and for canonicity (two nodes with the same sublevels are the same node).
\end{definition}

\begin{definition}[Well-formedness of an accumulated normal form]
\label{def:lvl-normlevel-wf}
\lean{Lean.Level.Normalize.NormLevel.WF, Lean.Level.Normalize.NormLevel.WF.of_mem, Lean.Level.Normalize.NormLevel.WF.sortedOf}
\leanok
\srcloc{Lean4Lean/Verify/Level.lean}{390}
\uses{def:lvl-extend, def:lvl-sorted}
\texttt{NormLevel.WF s} says that for every entry \texttt{(p, n)}: \texttt{p} is sorted; every
variable recorded in \texttt{n} has its name in \texttt{p}; and, when \texttt{p} is nonempty,
\texttt{p} extends by a single variable \texttt{v} some \texttt{p'} which is either the root or
another key, with \texttt{v} recorded at \texttt{p}. That last parent-pointer clause is what
later guarantees the sublevels are expressible as \texttt{imax} chains, and what licenses
\texttt{addConst} to drop $C(p,1)$ for nonempty \texttt{p}. \texttt{WF.of\_mem} climbs the parent
chain to produce, for any member of a key, a recorded witness at a subset of that key;
\texttt{WF.sortedOf} extracts sortedness of the root or of any key.
\end{definition}

\begin{theorem}[\texttt{normalizeAux} preserves well-formedness]
\label{thm:lvl-normalizeaux-wf}
\lean{Lean.Level.Normalize.normalizeAux_wf, Lean.Level.Normalize.normalizeAux_contains}
\leanok
\srcloc{Lean4Lean/Verify/Level.lean}{476}
\uses{def:lvl-normlevel-wf}
If the current path is the root or already a key of the accumulator, and the accumulator is well
formed, then \texttt{normalizeAux u path k acc} is well formed. \texttt{normalizeAux\_contains}
records separately that keys are never removed.
\end{theorem}
\begin{proof}
\uses{def:lvl-sorted, def:lvl-extend, def:lvl-normlevel-wf}
\leanok
Structural recursion following the twelve branches of \texttt{normalizeAux}, feeding
\texttt{addConst\_wf}, \texttt{addNode\_wf} and \texttt{addVar\_wf} with the
\texttt{Extend1}/\texttt{Sorted} facts supplied by \texttt{Sorted.orderedInsert}, and threading
the hypothesis through the recursive calls with \texttt{normalizeAux\_contains}.
\end{proof}

\begin{theorem}[\texttt{normalizeAux} is semantics-preserving]
\label{thm:lvl-normalizeaux-eval}
\lean{Lean.Level.Normalize.normalizeAux_eval}
\leanok
\srcloc{Lean4Lean/Verify/Level.lean}{575}
\uses{def:lvl-normlevel-eval, def:vlevel-oflevel}
If \texttt{VLevel.ofLevel ls u = some u'}, the current path is the root or a key, and the
accumulator is well formed, then
\texttt{(normalizeAux u path k acc).eval ls} $\rho$ is the maximum of
\texttt{acc.eval ls} $\rho$ and \texttt{evalPath path (u'.eval} $\rho$ \texttt{+ k)}. Each step
therefore adds exactly the sublevels of \texttt{u}, guarded by the current \texttt{imax}
conditions and shifted by the accumulated offset.
\end{theorem}
\begin{proof}
\uses{thm:lvl-normalizeaux-wf, fam:lvl-eval-le, def:lvl-normlevel-wf}
\leanok
Recursion on \texttt{normalizeAux}, one case per branch, about 85 lines. The interesting cases
are the \texttt{imax} distributions, discharged by the value-level identities \texttt{imax\_max}
and \texttt{imax\_imax}, and the two branches where the parameter is already in the path, where
\texttt{WF.of\_mem} produces the already-recorded witness that makes the new contribution
redundant.
\thesisref{axioms.tex, the \texttt{imax} rewriting rules of the inequality judgement; the same
identities, here at the level of values}
\end{proof}

\section{Subsumption and soundness of the comparison}

\begin{lemma}[The sorted-list subset test is a decision procedure]
\label{fam:lvl-subset}
\lean{Lean.Level.Normalize.subset_length, Lean.Level.Normalize.subset_mem, Lean.Level.Normalize.subset_eq, Lean.Level.Normalize.subset_of_sorted, Lean.Level.Normalize.subset_antisymm}
\leanok
\srcloc{Lean4Lean/Verify/Level.lean}{658}
The merge-style test \texttt{subset cmp l1 l2} implies that every element of \texttt{l1} lies in
\texttt{l2} and that \texttt{l1.length} $\le$ \texttt{l2.length}; on sorted lists the converse
holds as well, equal lengths force equality of the lists, and mutual inclusion is antisymmetric.
\end{lemma}
\begin{proof}
\uses{def:lvl-sorted}
\leanok
Each by induction on the second list following the three-way branch of \texttt{subset}, using
\texttt{LawfulBEqCmp.compare\_eq\_iff\_beq} in the \texttt{.eq} case; \texttt{subset\_of\_sorted}
additionally uses that the head of a sorted list is strictly below every later element.
\end{proof}

\begin{lemma}[Correctness of one subsumption step]
\label{fam:lvl-subsume-node}
\lean{Lean.Level.Normalize.subsumeVars_subset, Lean.Level.Normalize.subsumeVars_dominated, Lean.Level.Normalize.subsumeVars_eval, Lean.Level.Normalize.Node.subsumeBy_eval_le, Lean.Level.Normalize.Node.subsumeBy_eval_iff, Lean.Level.Normalize.Node.subsume_eval_le, Lean.Level.Normalize.Node.subsume_var_cases, Lean.Level.Normalize.Node.subsumeBy_const_drop}
\leanok
\srcloc{Lean4Lean/Verify/Level.lean}{718}
The lemmas describing \texttt{subsumeVars}, \texttt{Node.subsumeBy} and \texttt{Node.subsume}:
the result is always a sublist of the first argument; a variable that is dropped had a dominator
with the same name and a larger offset in the second argument; and, crucially,
\texttt{subsumeBy\_eval\_iff} says that if every variable of the dominating node has a nonzero
value and that node is already bounded by $m$, then \texttt{n1.subsumeBy same n2} is bounded by
$m$ exactly when \texttt{n1} is. Subsumption therefore never changes the total value.
\end{lemma}
\begin{proof}
\uses{fam:lvl-eval-le, fam:lvl-subset}
\leanok
Simultaneous induction on the two sorted variable lists along the three-way branch of
\texttt{subsumeVars}; the \texttt{eval\_iff} direction reads the bound off
\texttt{Node.eval\_le} and, for the constant, splits on \texttt{subsumeBy\_const\_drop}
(the constant is dropped only when it is below the dominating constant or below some dominating
variable). This is the paper's domination relation implemented as a merge over sorted lists; the
\texttt{same} flag, set when the two condition sets have equal length, is a genuine subtlety,
since at the same key a variable must not discharge itself.
\end{proof}

\begin{theorem}[Minimising an entry does not change the total]
\label{thm:lvl-minimize-eval}
\lean{Lean.Level.Normalize.NormLevel.minimize_eval_le, Lean.Level.Normalize.NormLevel.minimize_eval_iff}
\leanok
\srcloc{Lean4Lean/Verify/Level.lean}{933}
\texttt{minimize\_eval\_le} says that minimising a node only shrinks its value.
\texttt{minimize\_eval\_iff} is the converse under side conditions: if every variable of the map
belongs to its key, every entry other than \texttt{p1} is already bounded by $m$, and the
condition set \texttt{p1} is live under the valuation, then the minimised node is bounded by $m$
exactly when the original node is.
\end{theorem}
\begin{proof}
\uses{fam:lvl-subsume-node, fam:lvl-eval-le, fam:lvl-subset}
\leanok
Fold induction over the entry list of the map, applying \texttt{Node.subsumeBy\_eval\_iff} at
each step: an entry at a key below \texttt{p1} is live too, hence bounded by $m$, which is the
side condition that lemma needs.
\end{proof}

\begin{lemma}[One subsumption step changes only one key]
\label{thm:lvl-subsumption-step}
\lean{Lean.Level.Normalize.NormLevel.subsumption_step_get?}
\leanok
\srcloc{Lean4Lean/Verify/Level.lean}{1032}
The map obtained by replacing the entry at \texttt{p1} by its minimisation, or erasing it when
the minimisation is empty, agrees with the original at every other key, and at \texttt{p1} holds
the minimised node or nothing.
\end{lemma}
\begin{proof}
\leanok
Case split on emptiness and unfold \texttt{Std.TreeMap.getElem?\_erase} and
\texttt{getElem?\_insert}.
\end{proof}

\begin{definition}[Subsumption covers the original map]
\label{def:lvl-covers}
\lean{Lean.Level.Normalize.NormLevel.Covers, Lean.Level.Normalize.NormLevel.subsumption_covers, Lean.Level.Normalize.NormLevel.subsumption_vars}
\leanok
\srcloc{Lean4Lean/Verify/Level.lean}{1100}
\uses{thm:lvl-subsumption-step, fam:lvl-subsume-node, def:lvl-normlevel-wf}
\texttt{s'.Covers s} says that every variable recorded in \texttt{s} is still recorded in
\texttt{s'}, at a key all of whose elements lie in the original key.
\texttt{subsumption\_covers} proves that \texttt{s.subsumption} covers \texttt{s} and that each
of its entries is a sub-node of the corresponding entry of \texttt{s}; the proof is a fold
induction in which a dropped variable is replaced by the dominating witness at a strictly
smaller key, which later steps do not touch. \texttt{subsumption\_vars} preserves the
``variables belong to their key'' half of \texttt{WF}. \texttt{Covers} is exactly the part of a
map that the feasibility predicates of \S\ref{sec:lvl-reconstruction} read, which is why
feasibility survives subsumption.
\end{definition}

\begin{theorem}[Subsumption is semantics-preserving]
\label{thm:lvl-subsumption-eval}
\lean{Lean.Level.Normalize.NormLevel.subsumption_eval}
\leanok
\srcloc{Lean4Lean/Verify/Level.lean}{1157}
For a well-formed \texttt{s}, \texttt{s.subsumption} has the same value as \texttt{s} under every
valuation.
\end{theorem}
\begin{proof}
\uses{thm:lvl-minimize-eval, thm:lvl-subsumption-step, fam:lvl-eval-le, def:lvl-normlevel-wf}
\leanok
Fold induction over the entry list, maintaining that the accumulator has the same value as
\texttt{s}: each step uses \texttt{minimize\_eval\_iff} in one direction and
\texttt{minimize\_eval\_le} in the other, and observes that erasing a drained key has the same
effect on the value as keeping an empty node.
\end{proof}

\begin{theorem}[\texttt{normalize} is semantics-preserving]
\label{thm:lvl-normalize-eval}
\lean{Lean.Level.Normalize.normalize_eval}
\leanok
\srcloc{Lean4Lean/Verify/Level.lean}{1214}
\uses{def:lvl-normlevel-eval, def:vlevel-oflevel}
If \texttt{VLevel.ofLevel ls u = some u'} then \texttt{(normalize u).eval ls} $\rho$ equals
\texttt{u'.eval} $\rho$ for every valuation $\rho$. This is the soundness statement of the whole
accumulation phase.
\end{theorem}
\begin{proof}
\uses{thm:lvl-normalizeaux-eval, thm:lvl-subsumption-eval, thm:lvl-normalizeaux-wf}
\leanok
Combine \texttt{normalizeAux\_eval}, started from the empty and trivially well-formed map at the
root, with \texttt{subsumption\_eval}.
\thesisref{soundness.tex; the canonical form denotes the same function of the parameter values
as the input level}
\end{proof}

\begin{theorem}[Soundness of the normal-form comparison, Theorem 39]
\label{thm:lvl-le-eval}
\lean{Lean.Level.Normalize.NormLevel.le_eval}
\leanok
\srcloc{Lean4Lean/Verify/Level.lean}{1237}
If every variable recorded in \texttt{l2} belongs to its key and \texttt{l1.le l2} returns
\texttt{true}, then \texttt{l1.eval ls} $\rho \le$ \texttt{l2.eval ls} $\rho$ for every
valuation. Operationally, each entry of \texttt{l1} is folded against the entries of \texttt{l2},
each of which discharges what it can by \texttt{subsumeBy}, and the fold stopping with
\texttt{none} means nothing is left.
\end{theorem}
\begin{proof}
\uses{fam:lvl-subsume-node, fam:lvl-eval-le, fam:lvl-subset, family:treemap-axioms}
\leanok
Fix an entry of \texttt{l1} and a valuation making its condition set live. An entry of
\texttt{l2} at a key below it is then live too, hence bounded by the total of \texttt{l2}, which
is the side condition of \texttt{Node.subsumeBy\_eval\_iff}; induct along the entry list of
\texttt{l2} with that lemma at each discharging step. The docstring explains why a whole entry of
\texttt{l2} cannot be required to dominate a whole entry of \texttt{l1}, with the
\texttt{imax 2 v} $\le$ \texttt{max 2 v} example of \srcloc{Lean4Lean/Tests/Level.lean}{26}.
\thesisref{axioms.tex, the algorithmic level inequality; this is its semantic soundness
direction, following Théorème 39 of Géran's paper rather than the thesis rules}
\axfoot{Std.TreeMap.all\_eq\_all\_toList (Verify/Axioms.lean:10)}
\end{proof}

\begin{lemma}[\texttt{BEq}-equal normal forms have equal values]
\label{thm:lvl-eval-congr}
\lean{Lean.Level.Normalize.NormLevel.eval_congr, Lean.Level.Normalize.Node.eval_congr}
\leanok
\srcloc{Lean4Lean/Verify/Level.lean}{1273}
If \texttt{a == b} as \texttt{NormLevel}s, which by definition is mutual entry-wise inclusion,
then \texttt{a} and \texttt{b} have the same value under every valuation.
\end{lemma}
\begin{proof}
\uses{fam:lvl-eval-le, family:treemap-axioms}
\leanok
Prove the $\le$ direction for one-sided entry-wise inclusion by a double fold induction over the
two entry lists, then apply it twice. The lemma is needed because \texttt{NormLevel} is a
\texttt{Std.TreeMap}, and two maps with the same entries need not be the same tree.
\axfoot{Std.TreeMap.all\_eq\_all\_toList (Verify/Axioms.lean:10)}
\end{proof}

\section{Reconstruction: trees and admissible \texttt{imax} chains}
\label{sec:lvl-reconstruction}

A normal form must be turned back into a \texttt{Lean.Level}, which means choosing an
\texttt{imax} chain for each key. The subtlety, stated in the docstring at
\srcloc{Lean4Lean/Level.lean}{183}, is that a chain is not free: each edge itself contributes a
sublevel $V(S',v_i,0)$, so only \emph{admissible} orders may be used, and the choice must depend
on the sublevels alone if the output is to be canonical.

\begin{definition}[Semantics of the reconstruction tree]
\label{def:lvl-tree-eval}
\lean{Lean.Level.Normalize.Tree.eval, Lean.Level.Normalize.Tree.evalChild, Lean.Level.Normalize.Tree.eval_eq}
\leanok
\srcloc{Lean4Lean/Verify/Level.lean}{1306}
\uses{def:lvl-normlevel-eval}
A \texttt{Tree} node evaluates to the maximum of its own node value and the values of its
children, each child guarded by \texttt{Nat.imax} of the child's value and the value of the edge
label. This mirrors exactly the level \texttt{imax (reify t) (param a)} that \texttt{reify} emits
for an edge.
\end{definition}

\begin{theorem}[\texttt{reify} preserves the value of a tree]
\label{thm:lvl-tree-reify-eval}
\lean{Lean.Level.Normalize.Tree.reify_eval, Lean.Level.Normalize.Tree.reifyChild_eval, Lean.Level.Normalize.Tree.plainOffset?_eval, Lean.Level.Normalize.Tree.reifyChild_ge}
\leanok
\srcloc{Lean4Lean/Verify/Level.lean}{1404}
\uses{def:lvl-tree-eval}
For every tree \texttt{t}, evaluating the \texttt{Lean.Level} that \texttt{t.reify} emits at the
valuation \texttt{evalParam ls} $\rho$ gives \texttt{Tree.eval ls} $\rho$ \texttt{t}.
\end{theorem}
\begin{proof}
\leanok
Mutual induction with \texttt{reifyChild\_eval}, which states the same for a child list taken
together with the node constant. The delicate part is the \texttt{plainOffset?} optimisation,
where the \texttt{imax} guard on a child is dropped: \texttt{plainOffset?\_eval} shows the plain
form \texttt{a+k} and the guarded form \texttt{imax (a+k) a} differ only when the parameter is
$0$, where the node's own constant already covers the difference. The optimisation is motivated
in the docstring at \srcloc{Lean4Lean/Level.lean}{264}: without it \texttt{u+1} would reify to
\texttt{max 1 (imax (u+1) u)}.
\end{proof}

\begin{definition}[Subtrees at a path, and how \texttt{modify} changes them]
\label{def:lvl-tree-at}
\lean{Lean.Level.Normalize.Tree.At, Lean.Level.Normalize.Tree.At.append, Lean.Level.Normalize.Tree.At.suffix, Lean.Level.Normalize.Tree.At.nil_inv, Lean.Level.Normalize.Tree.At.append_inv, Lean.Level.Normalize.Tree.At_modify_self, Lean.Level.Normalize.Tree.At_modify_of_ne, Lean.Level.Normalize.Tree.At_modify_inv}
\leanok
\srcloc{Lean4Lean/Verify/Level.lean}{1449}
\texttt{Tree.At t p t'} says that \texttt{t'} is the subtree of \texttt{t} reached by following
the labels of \texttt{p}, innermost first, which is the order \texttt{Tree.modify} takes. The
accompanying lemmas invert paths (a path passes through all of its suffixes) and describe
\texttt{Tree.modify}: a modification is visible at its own path, invisible at incomparable
paths, and every subtree of the modified tree is either a subtree of the original or the
modified one. About 200 lines of list and tree plumbing, including the generic helpers
\texttt{modifyAt\_eq}, \texttt{mem\_modifyAt} and \texttt{suffix\_concat}.
\end{definition}

\begin{theorem}[A tree is bounded iff all its sublevels are]
\label{thm:lvl-tree-eval-le-iff}
\lean{Lean.Level.Normalize.Tree.eval_le_iff, Lean.Level.Normalize.Tree.eval_le_of, Lean.Level.Normalize.Tree.At.le, Lean.Level.Normalize.Tree.At.edge_le}
\leanok
\srcloc{Lean4Lean/Verify/Level.lean}{1595}
\uses{def:lvl-tree-at, def:lvl-tree-eval}
\texttt{Tree.eval ls} $\rho$ \texttt{t} $\le m$ holds exactly when (i) for every subtree at a
path \texttt{p}, \texttt{evalPath p} of that subtree's node value is $\le m$, and (ii) for every
edge \texttt{a :: p}, the sublevel the edge itself contributes, \texttt{evalPath (a :: p)} of the
value of \texttt{a}, is $\le m$.
\end{theorem}
\begin{proof}
\uses{fam:lvl-eval-le}
\leanok
The forward direction is \texttt{At.le} and \texttt{At.edge\_le}, that each subtree's and each
edge's contribution is below the total; the converse is the mutual
\texttt{eval\_le\_of}/\texttt{evalChild\_le\_of} induction. Clause (ii) is the crux of the whole
reconstruction argument: an \texttt{imax} chain is not free, so chains must be chosen admissibly.
\end{proof}

\begin{definition}[Dominated edges, feasible keys, admissible chains]
\label{def:lvl-dom-feas}
\lean{Lean.Level.Normalize.NormLevel.Dom, Lean.Level.Normalize.NormLevel.Feas, Lean.Level.Normalize.NormLevel.Adm, Lean.Level.Normalize.NormLevel.Dom.mono, Lean.Level.Normalize.NormLevel.Dom.le, Lean.Level.Normalize.NormLevel.Feas.mono, Lean.Level.Normalize.NormLevel.Feas.exchange, Lean.Level.Normalize.NormLevel.Feas.peel, Lean.Level.Normalize.NormLevel.Feas.perm, Lean.Level.Normalize.NormLevel.Feas.cons_last, Lean.Level.Normalize.NormLevel.Adm.suffix}
\leanok
\srcloc{Lean4Lean/Verify/Level.lean}{1616}
\uses{def:lvl-normlevel-eval, fam:lvl-eval-le}
\texttt{Dom s a acc} says that the sublevel contributed by an \texttt{imax} edge labelled
\texttt{a} is dominated by the map: some key \texttt{p} of \texttt{s} records a variable named
\texttt{a} and every element of \texttt{p} is \texttt{a} or lies in \texttt{acc}.
\texttt{Dom.le} turns this into the semantic statement that the edge contributes nothing beyond
the value of \texttt{s}, and \texttt{Dom.mono} says domination is monotone in \texttt{acc}, which
is what makes the greedy search of \texttt{lexChain} complete. \texttt{Feas s acc rem} says the
names in \texttt{rem} can be added to \texttt{acc} one at a time, each addition dominated; the
exchange lemma shows any dominated element may be taken first, and \texttt{peel} extracts the
element added last. \texttt{Adm l} says an already-ordered chain has every edge dominated.
\end{definition}

\begin{lemma}[The \texttt{addable} and \texttt{feasible} tests decide \texttt{Dom} and \texttt{Feas}]
\label{thm:lvl-addable-feasible}
\lean{Lean.Level.Normalize.NormLevel.addable_sound, Lean.Level.Normalize.NormLevel.addable_complete, Lean.Level.Normalize.NormLevel.feasible_sound, Lean.Level.Normalize.NormLevel.feasible_complete, Lean.Level.Normalize.NormLevel.feasible_go_sound, Lean.Level.Normalize.NormLevel.feasible_go_complete}
\leanok
\srcloc{Lean4Lean/Verify/Level.lean}{1639}
\texttt{s.addable a acc} decides \texttt{s.Dom a acc}, soundness unconditionally and completeness
assuming the keys of \texttt{s} and \texttt{acc} are sorted; \texttt{s.feasible acc rem} decides
\texttt{s.Feas acc rem} under the same sortedness assumptions. The fuel of \texttt{feasible.go}
is the length of \texttt{rem}, which the completeness proof shows suffices.
\end{lemma}
\begin{proof}
\uses{def:lvl-dom-feas, fam:lvl-subset, def:lvl-sorted, family:treemap-axioms}
\leanok
Soundness unfolds the \texttt{any}/\texttt{find?} and reads off the witness. Completeness turns
the membership statement back into a \texttt{subset} test with \texttt{subset\_of\_sorted}, using
that a sorted key is duplicate-free to erase the name being added; for \texttt{feasible} it also
uses \texttt{Feas.exchange} to justify committing to the first \texttt{addable} element.
\axfoot{Std.TreeMap.any\_eq\_any\_toList (Verify/Axioms.lean:14)}
\end{proof}

\begin{theorem}[Every key of a normal form admits an \texttt{imax} chain]
\label{thm:lvl-wf-feas}
\lean{Lean.Level.Normalize.NormLevel.WF.feas, Lean.Level.Normalize.NormLevel.Dom.mono_map, Lean.Level.Normalize.NormLevel.Feas.mono_map, Lean.Level.Normalize.normalize_feas}
\leanok
\srcloc{Lean4Lean/Verify/Level.lean}{1771}
\uses{def:lvl-normlevel-wf, def:lvl-dom-feas}
For a well-formed map, every key is \texttt{Feas} from the empty condition set; and since
subsumption \texttt{Covers} the map it was built from, and domination only reads variable names
and their keys, this survives into \texttt{normalize u} (\texttt{normalize\_feas}). Every key of
a normal form can therefore actually be reified as an \texttt{imax} chain.
\end{theorem}
\begin{proof}
\uses{def:lvl-covers}
\leanok
Induction on the length of the key, using the parent-pointer clause of \texttt{WF}: the parent is
the root or a key with one condition fewer, and the variable relating them dominates the edge
between them. \texttt{Feas.perm} and \texttt{Feas.cons\_last} reorder the chain so that the
variable just removed is the one added last. This is where the otherwise mysterious third clause
of \texttt{NormLevel.WF} pays off.
\end{proof}

\begin{theorem}[\texttt{lexChain} returns an admissible chain]
\label{thm:lvl-lexchain-spec}
\lean{Lean.Level.Normalize.NormLevel.lexChain_spec, Lean.Level.Normalize.NormLevel.lexChain_perm, Lean.Level.Normalize.NormLevel.lexChain_inj}
\leanok
\srcloc{Lean4Lean/Verify/Level.lean}{1836}
\uses{def:lvl-dom-feas, def:lvl-sorted}
\texttt{lexChain fuel p} always returns a permutation of \texttt{p}, including in the fallback
branch where it returns \texttt{p} itself (\texttt{lexChain\_perm}, no hypotheses). When the keys
of \texttt{s} are sorted, \texttt{fuel} is at least the length of \texttt{p}, \texttt{p} is
sorted and \texttt{p} is feasible, the result is in addition admissible. Distinct sorted keys
therefore get distinct chains (\texttt{lexChain\_inj}).
\end{theorem}
\begin{proof}
\uses{thm:lvl-addable-feasible, family:treemap-axioms}
\leanok
Recursion on the fuel. In the \texttt{find?} branch the element found is \texttt{addable} and
leaves a feasible remainder, so the recursive call applies; in the \texttt{none} branch
\texttt{Feas.peel} exhibits an element that \texttt{find?} would have found, a contradiction.
Injectivity follows from \texttt{lexChain\_perm} plus antisymmetry of permutation on sorted
lists. Taking the lexicographically least admissible chain is what makes reconstruction depend
only on the sublevels rather than on which \texttt{imax} chains happened to appear in the input,
the bug documented at \srcloc{Lean4Lean/Tests/Level.lean}{20}.
\axfoot{Std.TreeMap.any\_eq\_any\_toList (Verify/Axioms.lean:14)}
\end{proof}

\begin{theorem}[What \texttt{toTree} writes and what the tree contains]
\label{thm:lvl-totree-spec}
\lean{Lean.Level.Normalize.NormLevel.toTree_spec, Lean.Level.Normalize.NormLevel.WrittenAt, Lean.Level.Normalize.NormLevel.Accounted, Lean.Level.Normalize.NormLevel.treeVar, Lean.Level.Normalize.NormLevel.WrittenAt.write}
\leanok
\srcloc{Lean4Lean/Verify/Level.lean}{2137}
\uses{def:lvl-tree-at, thm:lvl-lexchain-spec}
After the single fold that builds the tree, every entry \texttt{(p, n)} is \texttt{WrittenAt} the
end of \texttt{p}'s chain, carrying \texttt{n}'s constant and \texttt{treeVar p n}, which is
\texttt{n}'s variables minus the one the innermost edge already contributes. Conversely the tree
is \texttt{Accounted} for: every subtree is either empty scaffolding or exactly such a written
entry, and every nonempty path is a suffix of some key's chain.
\end{theorem}
\begin{proof}
\uses{fam:lvl-subsume-node}
\leanok
Fold induction: a write is visible at its own path (\texttt{At\_modify\_self}) and leaves other
entries alone, either because the paths differ, by \texttt{lexChain\_inj} on distinct sorted
keys, or because the key repeats and the same data is written.
\end{proof}

\begin{theorem}[Reconstruction is semantics-preserving]
\label{thm:lvl-totree-eval}
\lean{Lean.Level.Normalize.NormLevel.toTree_eval, Lean.Level.Normalize.NormLevel.toTree_le_iff}
\leanok
\srcloc{Lean4Lean/Verify/Level.lean}{2217}
\uses{def:lvl-tree-eval, def:lvl-normlevel-eval}
For a map whose keys are sorted and all feasible,
\texttt{Tree.eval ls} $\rho$ \texttt{(toTree s)} equals \texttt{s.eval ls} $\rho$.
\texttt{toTree\_le\_iff} is the intermediate biconditional: the tree is bounded by $m$ exactly
when, for each entry, the recorded node and each edge of its chain are.
\end{theorem}
\begin{proof}
\uses{thm:lvl-totree-spec, thm:lvl-tree-eval-le-iff, thm:lvl-lexchain-spec, def:lvl-dom-feas, fam:lvl-subsume-node, family:treemap-axioms}
\leanok
Both inequalities come from \texttt{toTree\_le\_iff}. Downwards, the recorded node carries a
subset of the entry's sublevels and every edge is dominated, because \texttt{lexChain} emits only
admissible chains and \texttt{Dom.le} bounds a dominated edge. Upwards, the one sublevel the
recorded node omits, the innermost chain element at offset $0$, is exactly what the edge into
that node contributes.
\axfoot{Std.TreeMap.any\_eq\_any\_toList (Verify/Axioms.lean:14)}
\end{proof}

\section{Completeness and canonicity}

\begin{definition}[Sublevels and syntactic domination]
\label{def:lvl-sub}
\lean{Lean.Level.Normalize.Sub, Lean.Level.Normalize.Sub.path, Lean.Level.Normalize.Sub.le, Lean.Level.Normalize.Sub.le.trans, Lean.Level.Normalize.Sub.le.antisymm, Lean.Level.Normalize.NormLevel.HasSub, Lean.Level.Normalize.Sub.eval, Lean.Level.Normalize.NormLevel.HasSub.le_eval, Lean.Level.Normalize.Node.HasSub, Lean.Level.Normalize.NormLevel.hasSub_iff}
\leanok
\srcloc{Lean4Lean/Verify/Level.lean}{2535}
\uses{def:lvl-normlevel-eval, fam:lvl-subset}
A \texttt{Sub} is a single sublevel of the canonical form: \texttt{const p k} is $C(p,k)$ and
\texttt{var p x k} is $V(p,x,k)$. \texttt{Sub.le} is the paper's domination order: the
dominator's condition set must be a \emph{subset} of the dominated one, so that it fires
whenever the dominated sublevel does; a constant is dominated by a variable sublevel of offset
$K$ up to $K+1$; and a variable sublevel is dominated only by the same variable at a larger
offset. \texttt{HasSub s t} says \texttt{t} is recorded in \texttt{s}, and
\texttt{HasSub.le\_eval} that a recorded sublevel is below the total. Domination is transitive,
and antisymmetric on sorted condition sets.
\thesisref{axioms.tex, the level inequality judgement; Théorème 39 of Géran's paper is the
statement that this syntactic order decides the semantic one}
\end{definition}

\begin{theorem}[Separation: a semantic bound forces a syntactic dominator]
\label{thm:lvl-separation}
\lean{Lean.Level.Normalize.NormLevel.separation, Lean.Level.Normalize.NormLevel.bound, Lean.Level.Normalize.NormLevel.bound_spec, Lean.Level.Normalize.NormLevel.lt_eval, Lean.Level.Normalize.Node.lt_eval}
\leanok
\srcloc{Lean4Lean/Verify/Level.lean}{2674}
\uses{def:lvl-sub, def:lvl-normlevel-eval}
The converse of Theorem 39. If every name in every key of \texttt{l1} occurs in \texttt{ls},
every variable of \texttt{l1} belongs to its key, and \texttt{l1.eval ls} $\rho \le$
\texttt{l2.eval ls} $\rho$ for \emph{every} valuation $\rho$, then every sublevel of \texttt{l1}
has a syntactic dominator among the sublevels of \texttt{l2}.
\end{theorem}
\begin{proof}
\uses{fam:lvl-eval-le}
\leanok
Construct, for the sublevel under consideration, a separating valuation: every name in its
condition set is set to $1$, its own variable if it has one to a value $N$ exceeding every
constant and offset appearing in \texttt{l2} (\texttt{bound\_spec}), and everything else to $0$.
Only entries of \texttt{l2} at condition sets below the sublevel's can then contribute, and only
a sublevel with the same variable can reach $N$; reading off the entry that exceeds the bound
(\texttt{lt\_eval}) yields the dominator. About 75 lines, with no automation black boxes.
\thesisref{the converse direction of the algorithmic inequality of axioms.tex, which the thesis
states but does not prove}
\end{proof}

\begin{theorem}[Completeness of the normal-form comparison]
\label{thm:lvl-le-complete}
\lean{Lean.Level.Normalize.NormLevel.le_complete, Lean.Level.Normalize.NormLevel.le_fold_complete, Lean.Level.Normalize.subsumeVars_complete, Lean.Level.Normalize.Node.subsumeBy_var_complete, Lean.Level.Normalize.Node.subsumeBy_const_complete}
\leanok
\srcloc{Lean4Lean/Verify/Level.lean}{2886}
\uses{def:lvl-sub, def:lvl-varssorted}
If every sublevel of \texttt{l1} has a dominator among the sublevels of \texttt{l2}, and both
maps have sorted keys and sorted variable lists and \texttt{l1} has no empty entry, then
\texttt{l1.le l2} returns \texttt{true}. \texttt{le\_fold\_complete} is the fold-level statement:
an entry each of whose sublevels has a dominator in the remaining list is fully discharged.
\end{theorem}
\begin{proof}
\uses{fam:lvl-subsume-node, fam:lvl-subset, family:treemap-axioms}
\leanok
Per entry of \texttt{l1}, induct along the entry list of \texttt{l2}. Each step either skips a key
that is not a subset, or discharges through \texttt{subsumeBy}; exactness of \texttt{subsumeBy}
(\texttt{subsumeVars\_complete}, \texttt{subsumeBy\_var\_complete},
\texttt{subsumeBy\_const\_complete}) guarantees that a sublevel surviving a step still has a
dominator further down the list.
\axfoot{Std.TreeMap.all\_eq\_all\_toList (Verify/Axioms.lean:10)}
\end{proof}

\begin{theorem}[Minimisation removes exactly the dominated sublevels]
\label{thm:lvl-minimize-exact}
\lean{Lean.Level.Normalize.NormLevel.minimize_exact, Lean.Level.Normalize.NormLevel.minimize_exact_aux, Lean.Level.Normalize.NormLevel.minimize_hasSub, Lean.Level.Normalize.Node.subsume_hasSub, Lean.Level.Normalize.NormLevel.subsumption_step_hasSub}
\leanok
\srcloc{Lean4Lean/Verify/Level.lean}{3520}
\uses{def:lvl-sub, def:lvl-varssorted}
For a map with sorted keys and sorted variable lists, if \texttt{t} is a sublevel of the
minimised entry at \texttt{p1} and \texttt{t'} is any sublevel of the map with
\texttt{t.le t'}, then \texttt{t = t'}: minimisation leaves no sublevel that another sublevel
strictly dominates.
\end{theorem}
\begin{proof}
\uses{fam:lvl-subsume-node, fam:lvl-subset}
\leanok
Fold induction over the entry list (\texttt{minimize\_exact\_aux}, about 90 lines), splitting on
whether the dominating entry sits at the same key, where the equal-length \texttt{same} flag
applies and a variable must not discharge itself, or at a strictly smaller one.
\end{proof}

\begin{theorem}[Subsumption produces a reduced normal form]
\label{thm:lvl-subsumption-reduced}
\lean{Lean.Level.Normalize.NormLevel.Reduced, Lean.Level.Normalize.NormLevel.subsumption_reduced, Lean.Level.Normalize.normalize_reduced}
\leanok
\srcloc{Lean4Lean/Verify/Level.lean}{3536}
\uses{def:lvl-sub}
\texttt{Reduced s} says that domination between two recorded sublevels forces them to be equal.
\texttt{subsumption\_reduced} proves \texttt{s.subsumption} reduced for \texttt{s} with sorted
keys and sorted variable lists, and \texttt{normalize\_reduced} instantiates this at
\texttt{normalize u}.
\end{theorem}
\begin{proof}
\uses{thm:lvl-minimize-exact, thm:lvl-subsumption-step, def:lvl-varssorted}
\leanok
Fold induction maintaining that every already-processed key is minimised against the current
map; the key point is that later steps only shrink the map, so they cannot introduce new
domination.
\end{proof}

\begin{theorem}[Reduced forms with the same sublevels are equal]
\label{thm:lvl-eq-of-hassub}
\lean{Lean.Level.Normalize.NormLevel.eq_of_hasSub_iff, Lean.Level.Normalize.NormLevel.Reduced.hasSub_iff_hasSub, Lean.Level.Normalize.NormLevel.HasSub.path_sorted, Lean.Level.Normalize.sorted_pairs_eq}
\leanok
\srcloc{Lean4Lean/Verify/Level.lean}{3662}
\uses{def:lvl-sub, def:lvl-varssorted}
If two reduced normal forms with sorted keys mutually dominate each other's sublevels, then they
have literally the same sublevels; and two maps with the same sublevels, sorted variable lists
and no empty entries are \texttt{BEq}-equal.
\end{theorem}
\begin{proof}
\uses{thm:lvl-subsumption-reduced, fam:lvl-subset}
\leanok
Mutual domination plus reducedness plus antisymmetry of \texttt{Sub.le} on sorted condition sets
forces the dominating sublevel to be the sublevel itself. Equality of the maps then follows entry
by entry, using \texttt{VarsSorted.eq\_of\_mem\_iff} for the variable lists and
\texttt{sorted\_pairs\_eq} for the key lists.
\end{proof}

\begin{lemma}[Invariants of the output of \texttt{normalize}]
\label{fam:lvl-normalize-invariants}
\lean{Lean.Level.Normalize.normalize_vars_sorted, Lean.Level.Normalize.normalize_vars, Lean.Level.Normalize.normalize_sorted, Lean.Level.Normalize.normalize_sortedVars, Lean.Level.Normalize.normalize_nonempty, Lean.Level.Normalize.normalize_keys, Lean.Level.Normalize.normalize_feas, Lean.Level.Normalize.normalize_reduced}
\leanok
\srcloc{Lean4Lean/Verify/Level.lean}{1222}
The seven properties of \texttt{normalize u} that the end-to-end theorems take as separate
hypotheses: every recorded variable belongs to its key; keys are sorted; variable lists are
sorted by name; no entry is empty; every name in a key occurs in the input level, hence in
\texttt{ls} when the level translates; every key admits an admissible \texttt{imax} chain; and no
sublevel is dominated by another.
\end{lemma}
\begin{proof}
\uses{thm:lvl-normalizeaux-wf, def:lvl-covers, def:lvl-varssorted, thm:lvl-wf-feas, thm:lvl-subsumption-reduced}
\leanok
Each is the corresponding invariant of the accumulation phase, transported through subsumption
by \texttt{subsumption\_vars}, \texttt{subsumption\_covers}, \texttt{subsumption\_sortedVars} and
\texttt{subsumption\_reduced}; the statements live at lines 1222, 1807, 2413, 2443, 2524 and
3600. Collecting them as separate side conditions rather than one bundled predicate is what makes
the end-to-end theorems below take five or six hypotheses each.
\end{proof}

\begin{theorem}[Canonicity of the normal form]
\label{thm:lvl-normalize-complete}
\lean{Lean.Level.Normalize.normalize_complete}
\leanok
\srcloc{Lean4Lean/Verify/Level.lean}{3721}
\uses{def:vlevel-oflevel}
For levels that translate, \texttt{normalize u == normalize v} holds if and only if \texttt{u'}
and \texttt{v'} are equivalent as \texttt{VLevel}s, that is, have the same value under every
valuation.
\end{theorem}
\begin{proof}
\uses{thm:lvl-normalize-eval, thm:lvl-eval-congr, thm:lvl-separation, thm:lvl-eq-of-hassub, fam:lvl-normalize-invariants, family:treemap-axioms}
\leanok
Soundness is \texttt{normalize\_eval} followed by \texttt{eval\_congr}. For completeness, from
semantic equivalence each map bounds the other under every valuation, so \texttt{separation}
gives mutual domination of sublevels; both maps are reduced, so
\texttt{Reduced.hasSub\_iff\_hasSub} upgrades this to equality of the sublevel sets and
\texttt{eq\_of\_hasSub\_iff} to \texttt{BEq} equality.
\thesisref{axioms.tex, level equivalence defined as mutual inequality; this is a decision
procedure for it}
\axfoot{Std.TreeMap.all\_eq\_all\_toList (Verify/Axioms.lean:10)}
\end{proof}

\section{The flat fast path}

About 99.8\% of the levels arising in Lean, Std and Batteries contain no essential \texttt{imax}
(\srcloc{Lean4Lean/Level.lean}{302}), and for those the map can be skipped entirely. The fast
path is a performance optimisation whose transparency still has to be proved.

\begin{definition}[Flat normal forms]
\label{def:lvl-flat}
\lean{Lean.Level.Normalize.flatGet, Lean.Level.Normalize.NormLevel.Flat, Lean.Level.Normalize.RepAcc, Lean.Level.Normalize.NormLevel.addConst_flat, Lean.Level.Normalize.NormLevel.addNode_flat, Lean.Level.Normalize.toNormLevel_flat, Lean.Level.Normalize.NormLevel.subsumption_flat}
\leanok
\srcloc{Lean4Lean/Verify/Level.lean}{2948}
\uses{def:lvl-varssorted}
\texttt{flatGet c vs} describes, pointwise, the map that a level with no essential \texttt{imax}
normalises to: $C(\emptyset,c)$ at the root and $V(\{x\},x,k)$ at each singleton key.
\texttt{Flat s c vs} says \texttt{s} agrees with it at every key, and \texttt{RepAcc} relates the
two branches of the fast traversal to the accumulator of \texttt{normalizeAux}. The constructors
\texttt{addConst}, \texttt{addNode}, \texttt{toNormLevel} and \texttt{subsumption} all preserve
flatness, subsumption acting as the identity since nothing is ever dominated.
\end{definition}

\begin{lemma}[Everything downstream depends on the map only through its entry list]
\label{fam:lvl-congr}
\lean{Lean.Level.Normalize.toList_eq_of_get?_eq, Lean.Level.Normalize.normalizeAux_congr, Lean.Level.Normalize.NormLevel.subsumption_congr, Lean.Level.Normalize.NormLevel.toList_eq, Lean.Level.Normalize.NormLevel.addable_congr, Lean.Level.Normalize.NormLevel.feasible_congr, Lean.Level.Normalize.NormLevel.lexChain_congr, Lean.Level.Normalize.NormLevel.toTree_congr}
\leanok
\srcloc{Lean4Lean/Verify/Level.lean}{3177}
Two \texttt{Std.TreeMap}s with the same entries need not be the same tree, since the shape
depends on insertion order, so each consumer must be shown to factor through \texttt{toList}:
\texttt{normalizeAux}, \texttt{subsumption}, \texttt{addable}, \texttt{feasible},
\texttt{lexChain} and \texttt{toTree} all do, and \texttt{BEq}-equal maps have equal
\texttt{toList}s.
\end{lemma}
\begin{proof}
\uses{fam:lvl-subset, family:treemap-axioms}
\leanok
\texttt{toList\_eq\_of\_get?\_eq} compares the two sorted key-value lists pointwise
(\texttt{sorted\_pairs\_eq}); the rest propagate that equality through the definitions, one
congruence lemma per operation. A necessary but tedious layer, well flagged by the docstrings at
lines 3173--3176 and 3736--3739.
\axfoot{Std.TreeMap.all\_eq\_all\_toList, Std.TreeMap.any\_eq\_any\_toList (Verify/Axioms.lean:10,14)}
\end{proof}

\begin{theorem}[The flat traversal tracks \texttt{normalizeAux}]
\label{thm:lvl-flataux-rep}
\lean{Lean.Level.Normalize.flatAux_rep, Lean.Level.Normalize.NormLevel.Flat.toTree, Lean.Level.Normalize.NormLevel.Flat.toList, Lean.Level.Normalize.NormLevel.Flat.lexChain_singleton}
\leanok
\srcloc{Lean4Lean/Verify/Level.lean}{3261}
\uses{def:lvl-flat}
\texttt{flatAux l k (c, vs)} either returns flat data describing the same map that
\texttt{normalizeAux l [] k s} builds, or throws a map agreeing with it entry for entry. And a
flat map's \texttt{toTree} is exactly \texttt{flatTree c vs}: every key is a singleton, so its
\texttt{lexChain} is forced, and the fold adds one child per parameter in name order.
\end{theorem}
\begin{proof}
\uses{fam:lvl-congr, thm:lvl-lexchain-spec, family:treemap-axioms}
\leanok
Induction on the level for \texttt{flatAux\_rep}, with the \texttt{imax} case and the case of a
\texttt{max} whose first argument already threw handled by \texttt{normalizeAux\_congr};
\texttt{Flat.toTree} computes the entry list of a flat map and folds explicitly.
\axfoot{Std.TreeMap.any\_eq\_any\_toList (Verify/Axioms.lean:14)}
\end{proof}

\begin{theorem}[The fast path is transparent]
\label{thm:lvl-normalize-prime-eq}
\lean{Lean.Level.Normalize.normalize'_eq}
\leanok
\srcloc{Lean4Lean/Verify/Level.lean}{3793}
\uses{def:kernel-level-normalize}
\texttt{normalize' l} equals \texttt{(normalize l).toTree.reify} for every level: the flat fast
path computes exactly the tree the general path does.
\end{theorem}
\begin{proof}
\uses{thm:lvl-flataux-rep, fam:lvl-congr, def:lvl-flat, family:treemap-axioms}
\leanok
Case on \texttt{flatAux l 0 (0, [])}. In the \texttt{ok} case use \texttt{subsumption\_flat} and
\texttt{Flat.toTree}; in the \texttt{error} case use \texttt{subsumption\_congr} and
\texttt{toTree\_congr} to move from the map seeded by \texttt{toNormLevel} to the one the general
path builds.
\axfoot{Std.TreeMap.any\_eq\_any\_toList (Verify/Axioms.lean:14)}
\end{proof}

\section{End-to-end level theorems}

The seven results below are what the rest of the project could use. Only
\ref{thm:lvl-isequivlist-wf} actually reaches the verified type checker; see the review notes.
The fast-path disjuncts defer to core Lean's own \texttt{isEquiv}/\texttt{geq}, already verified
as \texttt{isEquiv\_wf} and \texttt{geq\_wf} in \srcloc{Lean4Lean/Verify/LevelStd.lean}{528}.

\begin{theorem}[Soundness of the level equivalence test]
\label{thm:lvl-isequiv-prime-wf}
\lean{Lean.Level.isEquiv'_wf}
\leanok
\srcloc{Lean4Lean/Verify/Level.lean}{3808}
\uses{def:kernel-level-isequiv, def:vlevel-oflevel}
If \texttt{isEquiv' u v} returns \texttt{true} and \texttt{u}, \texttt{v} translate to
\texttt{u'}, \texttt{v'}, then \texttt{u'} and \texttt{v'} are equivalent \texttt{VLevel}s.
\end{theorem}
\begin{proof}
\uses{thm:lvl-normalize-eval, thm:lvl-eval-congr, family:level-axioms, family:treemap-axioms}
\leanok
Split the disjunction defining \texttt{isEquiv'}: the fast path is core Lean's \texttt{isEquiv},
verified in \srcloc{Lean4Lean/Verify/LevelStd.lean}{528}; the fallback rewrites both sides with
\texttt{normalize\_eval} and applies \texttt{NormLevel.eval\_congr}.
\axfoot{Lean.Level.normalize\_eq, isExplicitSubsumedAux\_eq, instLawfulBEqLevel; Std.TreeMap.all\_eq\_all\_toList}
\end{proof}

\begin{theorem}[Soundness of the level-list equivalence test]
\label{thm:lvl-isequivlist-wf}
\lean{Lean.Level.isEquivList_wf}
\leanok
\srcloc{Lean4Lean/Verify/Level.lean}{3837}
\uses{def:kernel-level-isequiv}
If \texttt{Level.isEquivList us vs} holds and both lists translate elementwise, then the
translations are pointwise equivalent.
\end{theorem}
\begin{proof}
\uses{thm:lvl-isequiv-prime-wf, family:level-axioms, family:treemap-axioms}
\leanok
Induction on the two lists, applying \texttt{isEquiv'\_wf} pointwise. This is the only end-to-end
level theorem the verified type checker consumes
(\srcloc{Lean4Lean/Verify/TypeChecker/IsDefEq.lean}{398}), covering the constant-application
definitional-equality checks at \srcloc{Lean4Lean/TypeChecker.lean}{742} and
\srcloc{Lean4Lean/TypeChecker.lean}{866}.
\axfoot{Lean.Level.normalize\_eq, isExplicitSubsumedAux\_eq, instLawfulBEqLevel; Std.TreeMap.all\_eq\_all\_toList}
\end{proof}

\begin{theorem}[Soundness of the level inequality test]
\label{thm:lvl-geq-prime-wf}
\lean{Lean.Level.geq'_wf}
\leanok
\srcloc{Lean4Lean/Verify/Level.lean}{3828}
\uses{def:kernel-level-isequiv, def:vlevel-oflevel}
If \texttt{geq' u v} returns \texttt{true} and both levels translate, then \texttt{v'} $\le$
\texttt{u'} as \texttt{VLevel}s.
\end{theorem}
\begin{proof}
\uses{thm:lvl-le-eval, thm:lvl-normalize-eval, fam:lvl-normalize-invariants, family:level-axioms, family:treemap-axioms}
\leanok
Split the disjunction; the fast path is core's \texttt{geq}, verified as \texttt{geq\_wf} in
\srcloc{Lean4Lean/Verify/LevelStd.lean}{534}, and the fallback rewrites both sides with
\texttt{normalize\_eval} and applies \texttt{NormLevel.le\_eval} with the
\texttt{normalize\_vars} invariant.
\axfoot{Lean.Level.normalize\_eq, isExplicitSubsumedAux\_eq, instLawfulBEqLevel; Std.TreeMap.all\_eq\_all\_toList}
\end{proof}

\begin{theorem}[Soundness of reification]
\label{thm:lvl-normalize-prime-eval}
\lean{Lean.Level.normalize'_eval}
\leanok
\srcloc{Lean4Lean/Verify/Level.lean}{3822}
\uses{def:kernel-level-normalize, def:vlevel-oflevel}
For a level that translates, evaluating the reconstructed \texttt{Lean.Level}
\texttt{normalize' u} at \texttt{evalParam ls} $\rho$ gives \texttt{u'.eval} $\rho$: the
reconstruction denotes the same function as the input.
\end{theorem}
\begin{proof}
\uses{thm:lvl-normalize-prime-eq, thm:lvl-tree-reify-eval, thm:lvl-totree-eval, thm:lvl-normalize-eval, fam:lvl-normalize-invariants, family:treemap-axioms}
\leanok
Chain \texttt{normalize'\_eq}, \texttt{Tree.reify\_eval} and \texttt{toTree\_eval} (with the
sortedness and feasibility invariants of \texttt{normalize}), then \texttt{normalize\_eval}.
\axfoot{Std.TreeMap.any\_eq\_any\_toList (Verify/Axioms.lean:14)}
\end{proof}

\begin{theorem}[Completeness of the level equivalence test]
\label{thm:lvl-isequiv-prime-complete}
\lean{Lean.Level.isEquiv'_complete}
\leanok
\srcloc{Lean4Lean/Verify/Level.lean}{3861}
\uses{def:kernel-level-isequiv, def:vlevel-oflevel}
For levels that translate, \texttt{isEquiv' u v} returns \texttt{true} if and only if \texttt{u'}
and \texttt{v'} are equivalent. With soundness this makes \texttt{isEquiv'} a decision procedure
for level equivalence.
\end{theorem}
\begin{proof}
\uses{thm:lvl-normalize-complete, family:level-axioms, family:treemap-axioms}
\leanok
Rewrite the fallback disjunct with \texttt{normalize\_complete}; the fast-path disjunct is
subsumed by soundness of core's \texttt{isEquiv}.
\thesisref{decidability of level equivalence, the judgement of axioms.tex}
\axfoot{Lean.Level.normalize\_eq, isExplicitSubsumedAux\_eq, instLawfulBEqLevel; Std.TreeMap.all\_eq\_all\_toList}
\end{proof}

\begin{theorem}[Canonicity of the reconstructed level]
\label{thm:lvl-normalize-prime-complete}
\lean{Lean.Level.normalize'_complete}
\leanok
\srcloc{Lean4Lean/Verify/Level.lean}{3849}
\uses{def:kernel-level-normalize, def:vlevel-oflevel}
For levels that translate, \texttt{normalize' u = normalize' v} if and only if \texttt{u'} and
\texttt{v'} are equivalent: semantically equal levels reify to \emph{syntactically} equal levels.
\end{theorem}
\begin{proof}
\uses{thm:lvl-normalize-complete, thm:lvl-normalize-prime-eval, fam:lvl-congr, thm:lvl-normalize-prime-eq, family:treemap-axioms}
\leanok
Left to right by \texttt{normalize'\_eval} at a fixed valuation. Right to left from
\texttt{normalize\_complete}, which gives \texttt{BEq}-equal normal forms, hence equal entry
lists, and then \texttt{toTree\_congr}. This is the strongest statement in the file, a genuine
canonical form for universe levels, which the thesis does not have.
\axfoot{Std.TreeMap.all\_eq\_all\_toList, Std.TreeMap.any\_eq\_any\_toList (Verify/Axioms.lean:10,14)}
\end{proof}

\begin{theorem}[Completeness of the level inequality test]
\label{thm:lvl-geq-prime-complete}
\lean{Lean.Level.geq'_complete}
\leanok
\srcloc{Lean4Lean/Verify/Level.lean}{3870}
\uses{def:kernel-level-isequiv, def:vlevel-oflevel}
For levels that translate, \texttt{geq' u v} returns \texttt{true} if and only if \texttt{v'}
$\le$ \texttt{u'}.
\end{theorem}
\begin{proof}
\uses{thm:lvl-geq-prime-wf, thm:lvl-separation, thm:lvl-le-complete, thm:lvl-normalize-eval, fam:lvl-normalize-invariants, family:level-axioms, family:treemap-axioms}
\leanok
Soundness is \texttt{geq'\_wf}. For completeness, \texttt{separation} applied to
\texttt{normalize v} and \texttt{normalize u} produces a dominator for each sublevel, and
\texttt{le\_complete} turns that into a \texttt{true} answer; the \texttt{refine} takes six of
the \texttt{normalize} invariants at once.
\thesisref{decidability of the algorithmic inequality of axioms.tex}
\axfoot{Lean.Level.normalize\_eq, isExplicitSubsumedAux\_eq, instLawfulBEqLevel; Std.TreeMap.all\_eq\_all\_toList}
\end{proof}

\section{The order used to sort \texttt{max} arguments}

Core Lean's \texttt{Level.normalize} sorts the arguments of a \texttt{max} with
\texttt{Array.qsort normLt}, so the verification of that function in
\srcloc{Lean4Lean/Verify/LevelStd.lean}{228} needs to know that \texttt{normLt} is a strict weak
order. \srcloc{Lean4Lean/Verify/NormLt.lean}{1} supplies exactly the two facts \texttt{qsort}
requires, by reformulating \texttt{normLt} as an \texttt{Ordering}-valued comparison.

\begin{definition}[Structural comparison of level bases]
\label{def:normlt-size-basecmp}
\lean{Lean.Level.baseCmp, Lean.Level.size, Lean.Level.structCmp}
\leanok
\srcloc{Lean4Lean/Verify/NormLt.lean}{46}
\texttt{baseCmp} compares two levels that are not \texttt{succ}s: \texttt{max} against
\texttt{max} and \texttt{imax} against \texttt{imax} by recursively comparing each argument's
base and then its offset, \texttt{param} against \texttt{param} and \texttt{mvar} against
\texttt{mvar} by \texttt{Name.cmp}, and levels with different constructors by the constructor tag
\texttt{ctorToNat}. It is defined by well-founded recursion on the sum of the structural sizes
\texttt{size} of the two levels, using that a level's \texttt{getLevelOffset} is no larger than
the level itself. \texttt{structCmp} is the auxiliary same-constructor part, returning
\texttt{.eq} on differing constructors. \texttt{size}, \texttt{structCmp} and the size lemmas are
private.
\end{definition}

\begin{definition}[Ordering-valued comparison of levels]
\label{def:normlt-normcmp}
\lean{Lean.Level.normCmp}
\leanok
\srcloc{Lean4Lean/Verify/NormLt.lean}{65}
\uses{def:normlt-size-basecmp}
\texttt{normCmp l1 l2} compares \texttt{l1} and \texttt{l2} by their bases, with
\texttt{baseCmp} on \texttt{getLevelOffset}, and then, on equal bases, by their \texttt{succ}
offsets. This is the \texttt{Ordering}-valued reformulation of core Lean's boolean
\texttt{Level.normLt}.
\end{definition}

\begin{lemma}[\texttt{baseCmp} is the lexicographic product of tag and structure]
\label{thm:normlt-basecmp-eq}
\lean{Lean.Level.baseCmp_eq}
\leanok
\srcloc{Lean4Lean/Verify/NormLt.lean}{78}
For all levels, \texttt{baseCmp l1 l2} is \texttt{(compare l1.ctorToNat l2.ctorToNat).then
(structCmp l1 l2)}: the constructor tags decide first and the structural comparison only breaks
ties.
\end{lemma}
\begin{proof}
\uses{def:normlt-size-basecmp}
\leanok
Case split on both levels, then \texttt{simp [baseCmp, structCmp, normCmp, ctorToNat]}. This is
the reformulation that lets the lexicographic machinery of \texttt{Lean4Lean.Std.Ord} apply.
\end{proof}

\begin{lemma}[\texttt{normCmp} is a lawful comparison]
\label{fam:normlt-order-instances}
\lean{Lean.Level.normCmp_swap, Lean.Level.normCmp_rot, Lean.Level.normCmp_refl, Lean.Level.eq_of_normCmp_eq, Lean.Level.baseCmp_swap, Lean.Level.baseCmp_rot, Lean.Level.baseCmp_refl}
\leanok
\srcloc{Lean4Lean/Verify/NormLt.lean}{99}
\uses{def:normlt-normcmp}
The family establishing that \texttt{normCmp}, and with it \texttt{baseCmp}, is an oriented,
transitive, reflexive comparison that decides equality: \texttt{normCmp l2 l1} is the swap of
\texttt{normCmp l1 l2}; the rotation property \texttt{Rot} holds at every triple, whence
transitivity; \texttt{normCmp l l} is \texttt{.eq}; and \texttt{normCmp l1 l2 = .eq} implies
\texttt{l1 = l2}. Five anonymous instances are derived from them
(\texttt{TransCmp} and \texttt{ReflCmp} for both comparisons, and \texttt{LawfulEqCmp normCmp}),
so their generated names are not stated here.
\end{lemma}
\begin{proof}
\uses{thm:normlt-basecmp-eq}
\leanok
Transitivity goes through \texttt{Rot.then'} and \texttt{Rot.of\_transCmp} applied to the
lexicographic decomposition \texttt{baseCmp\_eq}, with a strong induction on the sum of the sizes
of the three levels; this is the device shared with \texttt{Name.cmp}
(\srcloc{Lean4Lean/Std/Ord.lean}{25}), needed because \texttt{structCmp} returns \texttt{.eq} on
unrelated constructors and so is not itself transitive. Antisymmetry uses the private
\texttt{level\_ext}, that a level is determined by its base and its offset.
\end{proof}

\begin{lemma}[\texttt{normLtAux} accumulates offsets then runs \texttt{normCmp}]
\label{thm:normlt-aux-eq}
\lean{Lean.Level.normLtAux_eq}
\leanok
\srcloc{Lean4Lean/Verify/NormLt.lean}{264}
For all \texttt{l1}, \texttt{k1}, \texttt{l2}, \texttt{k2}, core Lean's
\texttt{normLtAux l1 k1 l2 k2} equals the boolean test that
\texttt{(baseCmp l1.getLevelOffset l2.getLevelOffset).then (compare (l1.getOffset + k1)
(l2.getOffset + k2))} is \texttt{.lt}. This is the bridge between the imperative-looking core
function and the algebraic comparison.
\end{lemma}
\begin{proof}
\uses{fam:normlt-order-instances, def:normlt-normcmp}
\leanok
Induction along \texttt{normLtAux.induct} with fourteen cases, one per syntactic branch of
\texttt{normLtAux}: the \texttt{succ} cases move a unit from the level into the accumulator, the
syntactically equal cases close by \texttt{baseCmp\_refl}, and the remaining cases case-split the
resulting \texttt{Ordering}. About 80 lines, the bulk of the file.
\end{proof}

\begin{theorem}[\texttt{normLt} is \texttt{normCmp = lt}]
\label{thm:normlt-eq}
\lean{Lean.Level.normLt_eq}
\leanok
\srcloc{Lean4Lean/Verify/NormLt.lean}{345}
\uses{def:normlt-normcmp}
For all levels, \texttt{normLt l1 l2} equals \texttt{normCmp l1 l2 == .lt}.
\end{theorem}
\begin{proof}
\uses{thm:normlt-aux-eq}
\leanok
Unfold \texttt{normLt} and apply \texttt{normLtAux\_eq} with both accumulators zero.
\end{proof}

\begin{lemma}[\texttt{normLt} is a strict weak order]
\label{thm:normlt-swo}
\lean{Lean.Level.normLt_asymm, Lean.Level.normLt_le_trans, Lean.Level.normLt_same_base}
\leanok
\srcloc{Lean4Lean/Verify/NormLt.lean}{350}
The three exported facts: \texttt{normLt a b} implies \texttt{¬ normLt b a}; \texttt{¬ normLt b
a} and \texttt{¬ normLt c b} imply \texttt{¬ normLt c a}; and on levels with equal bases,
\texttt{normLt} is the comparison of the offsets. The first two are exactly the hypotheses
\texttt{Array.qsort\_sorted} (\ref{thm:qsort-sorted}) takes.
\end{lemma}
\begin{proof}
\uses{thm:normlt-eq, fam:normlt-order-instances}
\leanok
Rewrite with \texttt{normLt\_eq} and apply \texttt{OrientedCmp.not\_lt\_of\_lt} and
\texttt{TransCmp.isGE\_trans} for the \texttt{normCmp} instances; \texttt{normLt\_same\_base}
rewrites with \texttt{baseCmp\_refl}. These three are the only results of the file that the rest
of the project consumes (\srcloc{Lean4Lean/Verify/LevelStd.lean}{240}); everything else is
private or an instance.
\end{proof}

\section{A specification for \texttt{Array.qsort}}

\srcloc{Lean4Lean/Verify/QSort.lean}{1} is vendored from the upstream Lean verification in
\texttt{leanprover/lean4\#14658} and carries a Lean FRO copyright naming Kim Morrison as author.
It exists only because \texttt{Array.qsort} has no specification in the standard library; the
consumer is \srcloc{Lean4Lean/Verify/LevelStd.lean}{228}. It is the only file in the project
using the Lean 4 module system, and it is \texttt{grind}-driven throughout.

\begin{lemma}[Permutations restricted to a slice]
\label{fam:qsort-perm-extract}
\lean{List.Perm.take_of_getElem, List.Perm.drop_of_getElem, Array.Perm.extract', Vector.Perm.extract'}
\leanok
\srcloc{Lean4Lean/Verify/QSort.lean}{58}
If \texttt{l1} is a permutation of \texttt{l2} and the two lists agree elementwise from index
\texttt{i} on (respectively below \texttt{i}), then their prefixes (respectively suffixes) of
length \texttt{i} are permutations; and correspondingly, if \texttt{xs} and \texttt{ys} are
permutations as arrays or vectors and agree outside the range \texttt{[lo, hi)}, then
\texttt{xs.extract lo hi} and \texttt{ys.extract lo hi} are permutations. These are what let the
partition step be composed with the recursive sorts.
\end{lemma}
\begin{proof}
\leanok
Reduce to the existing \texttt{List.Perm.take\_of\_getElem?} and \texttt{drop\_of\_getElem?}, then
to the array and vector versions through \texttt{perm\_iff\_toList\_perm}; the index bookkeeping
is closed by \texttt{grind}. All four are declared at \texttt{\_root\_}, so they land in the
\texttt{List}, \texttt{Array} and \texttt{Vector} namespaces, and are registered as global
\texttt{grind} lemmas at lines 105--106.
\end{proof}

\begin{theorem}[\texttt{qsort} preserves size]
\label{thm:qsort-size}
\lean{Array.size_qsort}
\leanok
\srcloc{Lean4Lean/Verify/QSort.lean}{110}
\texttt{(qsort as lt lo hi).size = as.size}.
\end{theorem}
\begin{proof}
\leanok
\texttt{grind [qsort]}.
\end{proof}

\begin{theorem}[\texttt{qsort} returns a permutation]
\label{thm:qsort-perm}
\lean{Array.qsort_perm}
\leanok
\srcloc{Lean4Lean/Verify/QSort.lean}{137}
\texttt{qsort as lt lo hi} is a permutation of \texttt{as}.
\end{theorem}
\begin{proof}
\uses{fam:qsort-perm-extract}
\leanok
From the private \texttt{qsort\_sort\_perm}, an induction along \texttt{qsort.sort}, and
\texttt{qpartition\_perm}, which in turn follows from \texttt{qpartition\_loop\_perm} and
\texttt{Vector.swap\_perm}; the top level is \texttt{grind [qsort]}. Consumed by
\srcloc{Lean4Lean/Verify/LevelStd.lean}{228} to justify that core \texttt{Level.normalize} may
sort the arguments of a \texttt{max}.
\end{proof}

\begin{lemma}[Partition specification and frame conditions]
\label{fam:qsort-partition-spec}
\lean{Array.qpartition_fst_lt_hi, Array.getElem_qsort_sort_of_lt_lo, Array.getElem_qsort_sort_of_hi_lt, Array.getElem_qsort_sort_mem, Array.qpartition_spec₁, Array.qpartition_spec₂}
\leanok
\srcloc{Lean4Lean/Verify/QSort.lean}{184}
The private core lemmas about one partition step: every element of the active range strictly
below the returned pivot index is \texttt{lt} the pivot, and every element strictly above it and
at most \texttt{hi} is not; when \texttt{lo < hi} and \texttt{lt} is asymmetric the returned pivot
index is $<$ \texttt{hi}; sorting a range leaves the entries outside it untouched; and every
entry the sort places in the range came from that range.
\end{lemma}
\begin{proof}
\uses{fam:qsort-perm-extract}
\leanok
Each is a \texttt{fun\_induction} on \texttt{qpartition.loop} or on \texttt{qsort.sort} closed by
\texttt{grind}, with the pivot-position lemma split into the two cases
\texttt{qpartition\_loop\_lt\_hi₁} and \texttt{qpartition\_loop\_lt\_hi₂}.
\end{proof}

\begin{theorem}[Adjacent elements of a sorted range are non-decreasing]
\label{thm:qsort-sort-spec}
\lean{Array.qsort_sort_spec}
\leanok
\srcloc{Lean4Lean/Verify/QSort.lean}{264}
Assume \texttt{lt} is asymmetric and its negation transitive. For
\texttt{as' = qsort.sort lt as lo hi} and \texttt{lo} $\le$ \texttt{i} $<$ \texttt{hi}, the pair
\texttt{as'[i]}, \texttt{as'[i+1]} is non-decreasing: \texttt{¬ lt as'[i+1] as'[i]}.
\end{theorem}
\begin{proof}
\uses{fam:qsort-partition-spec, fam:qsort-perm-extract}
\leanok
Strong recursion on \texttt{qsort.sort}. After partitioning at \texttt{mid}, the case
\texttt{i < mid} recurses on the lower half, using the frame lemma
\texttt{getElem\_qsort\_sort\_of\_lt\_lo} to see the upper half as unchanged; the case
\texttt{mid = i} is the interesting one, locating both elements as members of the respective
extracted ranges (\texttt{getElem\_qsort\_sort\_mem}) and concluding from the two partition specs
plus transitivity; the case \texttt{mid < i} recurses on the upper half. The longest proof in the
file, and the only place where both hypotheses on \texttt{lt} are used.
\end{proof}

\begin{lemma}[\texttt{qsort} output is adjacent-sorted]
\label{thm:qsort-sorted-adjacent}
\lean{Array.qsort_sorted₁', Array.qsort_sorted₁}
\leanok
\srcloc{Lean4Lean/Verify/QSort.lean}{348}
The public wrappers of \texttt{qsort\_sort\_spec}: for asymmetric \texttt{lt} with transitive
negation, adjacent elements of the sorted slice \texttt{[lo, hi]} of \texttt{as.qsort lt lo hi},
and of the whole array under \texttt{as.qsort lt}, are non-decreasing.
\end{lemma}
\begin{proof}
\uses{thm:qsort-sort-spec}
\leanok
Unfold \texttt{qsort}, dispatch the degenerate empty-array branch by \texttt{grind}, and apply
\texttt{qsort\_sort\_spec}.
\end{proof}

\begin{theorem}[\texttt{qsort} output is sorted at arbitrary index pairs]
\label{thm:qsort-sorted}
\lean{Array.qsort_sorted', Array.qsort_sorted}
\leanok
\srcloc{Lean4Lean/Verify/QSort.lean}{370}
For asymmetric \texttt{lt} with transitive negation: if \texttt{lo} $\le$ \texttt{i} $<$
\texttt{j} $\le$ \texttt{hi} then \texttt{¬ lt (as.qsort lt lo hi)[j] (as.qsort lt lo hi)[i]};
and in the whole-array variant, \texttt{i < j} implies
\texttt{¬ lt (as.qsort lt)[j] (as.qsort lt)[i]}.
\end{theorem}
\begin{proof}
\uses{thm:qsort-sorted-adjacent}
\leanok
Induction on \texttt{j}, chaining adjacent-sortedness with transitivity of the negation.
\texttt{Array.qsort\_sorted} is the theorem consumed by
\srcloc{Lean4Lean/Verify/LevelStd.lean}{241}.
\end{proof}

\section{The equivalence manager}

The kernel memoises structural equality of expressions in a union-find keyed by \texttt{ExprMap}
lookup, that is, by \texttt{Expr}'s \texttt{BEq}, which ignores binder names, \texttt{mdata}
payloads and \texttt{proj} structure names (\ref{def:kernel-equiv-manager}).
\srcloc{Lean4Lean/Verify/EquivManager.lean}{1} names that quotient, proves it semantically
harmless, and shows the cache sound. Two notions it uses are defined elsewhere and have no node
of their own: the congruence closure \texttt{IsDefEqE} that the manager is allowed to assert
(\srcloc{Lean4Lean/Verify/TypeChecker/Basic.lean}{49}) and the manager invariant
\texttt{EquivManager.WF} (\srcloc{Lean4Lean/Verify/TypeChecker/Basic.lean}{87}), which says that
the union-find's node ids are exactly the ids assigned to expressions, and that any two
identified ids belong to expressions related by \texttt{IsDefEqE}.

\begin{definition}[Relevant structural equality of expressions]
\label{ind:eqvmgr-relevanteq}
\lean{Lean4Lean.EquivManager.RelevantEq}
\leanok
\srcloc{Lean4Lean/Verify/EquivManager.lean}{10}
An inductively defined congruence on \texttt{Lean.Expr} that ignores exactly the parts the kernel
treats as irrelevant: binder names and binder info in \texttt{lam} and \texttt{forallE}, the
\texttt{mdata} payload, the non-dependent flag of \texttt{letE}, and the structure name in
\texttt{proj}, where only the projection index must agree. All other components must be
syntactically equal. This is the quotient the memoised \texttt{==} on \texttt{Expr} works up to.
\end{definition}

\begin{lemma}[\texttt{RelevantEq} is an equivalence refined by \texttt{BEq}]
\label{fam:eqvmgr-relevanteq-lemmas}
\lean{Lean4Lean.EquivManager.RelevantEq.rfl, Lean4Lean.EquivManager.RelevantEq.symm, Lean4Lean.EquivManager.RelevantEq.of_eqv, Lean4Lean.EquivManager.RelevantEq.trans}
\leanok
\srcloc{Lean4Lean/Verify/EquivManager.lean}{26}
\texttt{RelevantEq} is reflexive, symmetric and transitive, and \texttt{e1 == e2} implies
\texttt{RelevantEq e1 e2}.
\end{lemma}
\begin{proof}
\uses{ind:eqvmgr-relevanteq, family:expr-beq}
\leanok
Each by structural induction on the expression or on the \texttt{RelevantEq} derivation;
\texttt{of\_eqv} unfolds \texttt{Expr.eqv'} and closes the cases with \texttt{grind [RelevantEq]}.
Only \texttt{rfl} and \texttt{of\_eqv} are used anywhere; see the review notes.
\end{proof}

\begin{theorem}[\texttt{RelevantEq} implies \texttt{IsDefEqE}]
\label{thm:eqvmgr-isdefeqe-relevant}
\lean{Lean4Lean.EquivManager.IsDefEqE.relevant}
\leanok
\srcloc{Lean4Lean/Verify/EquivManager.lean}{70}
\uses{ind:eqvmgr-relevanteq}
If \texttt{RelevantEq e1 e2} then \texttt{IsDefEqE env us Δ e1 e2}: relevant structural equality
is a special case of the congruence closure the equivalence manager may assert.
\end{theorem}
\begin{proof}
\leanok
Induction on the \texttt{RelevantEq} derivation, mapping each constructor to the corresponding
\texttt{IsDefEqE} constructor.
\end{proof}

\begin{theorem}[Relevantly equal expressions translate to definitionally equal terms]
\label{thm:eqvmgr-relevanteq-uniq}
\lean{Lean4Lean.EquivManager.RelevantEq.uniq}
\leanok
\srcloc{Lean4Lean/Verify/EquivManager.lean}{81}
\uses{ind:eqvmgr-relevanteq, ind:trexprs, ind:vt-vlctx-isdefeq, def:hastype-istype}
Let \texttt{env} be a well-formed environment and let two local contexts be definitionally equal.
If \texttt{RelevantEq e1 e2}, and \texttt{e1} translates to \texttt{e1'} in the first context and
\texttt{e2} to \texttt{e2'} in the second (both by \texttt{TrExprS}), then \texttt{e1'} and
\texttt{e2'} are definitionally equal terms.
\end{theorem}
\begin{proof}
\uses{thm:trproj-uniq}
\leanok
Induction on the first translation derivation, inverting the \texttt{RelevantEq} and the second
translation in lockstep. The binder cases extend the context equality with
\texttt{VLCtx.IsDefEq.cons}; the \texttt{proj} case defers to \texttt{TrProj.uniq}
(\ref{thm:trproj-uniq}), which is stated but not proved, and which is the sole reason this
theorem and the two below carry \texttt{sorryAx}. That case is exactly where the semantic content
lies: ignoring the structure name of a projection is sound only if the translation of a
projection is determined by its index, which is what \texttt{TrProj.uniq} would say.
\axfoot{sorryAx via Lean4Lean.TrProj.uniq (Verify/Typing/Lemmas.lean:992)}
\end{proof}

\begin{theorem}[\texttt{IsDefEqE} transports translations]
\label{thm:eqvmgr-isdefeqe-trexpr}
\lean{Lean4Lean.EquivManager.IsDefEqE.trExpr}
\leanok
\srcloc{Lean4Lean/Verify/EquivManager.lean}{119}
\uses{def:trexpr, ind:trexprs}
If \texttt{IsDefEqE env Us Δ r1 r2} holds in a context with no bound variables, and \texttt{r1} is
relevantly equal to some \texttt{e1} that translates to \texttt{e'} in a weakened context
\texttt{Δ'}, then there is an \texttt{e2} relevantly equal to \texttt{r2} which also translates,
up to definitional equality, to \texttt{e'}.
\end{theorem}
\begin{proof}
\uses{thm:eqvmgr-relevanteq-uniq}
\leanok
Induction on the \texttt{IsDefEqE} derivation. The \texttt{defeq} case weakens both translations
along the bound-variable lift, using that the context has no bound variables to discharge the
lifting of loose indices, and applies \texttt{RelevantEq.uniq}; the congruence cases rebuild the
translation with the corresponding \texttt{TrExpr} congruence lemmas.
\axfoot{sorryAx via Lean4Lean.TrProj.uniq (Verify/Typing/Lemmas.lean:992)}
\end{proof}

\begin{theorem}[\texttt{IsDefEqE} is sound]
\label{thm:eqvmgr-isdefeqe-uniq}
\lean{Lean4Lean.EquivManager.IsDefEqE.uniq}
\leanok
\srcloc{Lean4Lean/Verify/EquivManager.lean}{163}
\uses{ind:trexprs, def:hastype-istype}
If \texttt{env} is well formed, the local context has no bound variables and is well formed, and
\texttt{e1}, \texttt{e2} translate to \texttt{e1'}, \texttt{e2'}, then
\texttt{IsDefEqE env Us Δ e1 e2} implies that \texttt{e1'} and \texttt{e2'} are definitionally
equal.
\end{theorem}
\begin{proof}
\uses{thm:eqvmgr-isdefeqe-trexpr, thm:eqvmgr-relevanteq-uniq}
\leanok
Apply \texttt{IsDefEqE.trExpr} at the identity weakening and reflexive \texttt{RelevantEq}, then
combine with \texttt{RelevantEq.uniq} and transitivity of definitional equality. This is the
theorem the type checker consumes (\srcloc{Lean4Lean/Verify/TypeChecker/IsDefEq.lean}{165}) to
turn a positive answer of the cached \texttt{isEquiv} into a real definitional equality.
\axfoot{sorryAx via Lean4Lean.TrProj.uniq (Verify/Typing/Lemmas.lean:992)}
\end{proof}

\begin{definition}[Monotone growth of the equivalence manager]
\label{def:eqvmgr-le}
\lean{Lean4Lean.EquivManager.LE, Lean4Lean.EquivManager.LE.rfl, Lean4Lean.EquivManager.LE.trans}
\leanok
\srcloc{Lean4Lean/Verify/EquivManager.lean}{169}
\uses{def:kernel-equiv-manager}
\texttt{m1} $\le$ \texttt{m2} records that the union-find of \texttt{m2} identifies everything
\texttt{m1} did, and that every expression already assigned a node id in \texttt{m1} keeps that
id in \texttt{m2}. It is reflexive and transitive, with an anonymous \texttt{LE EquivManager}
instance. The state only ever grows; this order is what lets the monadic reasoning below carry
facts forward across \texttt{bind}.
\end{definition}

\begin{definition}[A Hoare-style predicate for the state monad]
\label{def:eqvmgr-mwf}
\lean{Lean4Lean.EquivManager.M.WF, Lean4Lean.EquivManager.M.WF.stateWF, Lean4Lean.EquivManager.M.WF.pure, Lean4Lean.EquivManager.M.WF.bind, Lean4Lean.EquivManager.M.WF.mono, Lean4Lean.EquivManager.M.WF.bind_le, Lean4Lean.EquivManager.M.WF.andM}
\leanok
\srcloc{Lean4Lean/Verify/EquivManager.lean}{181}
\uses{def:eqvmgr-le}
\texttt{M.WF env Us Δ m x P} says: if the manager \texttt{m} is well formed, then running the
stateful computation \texttt{x} from \texttt{m} yields a still well-formed manager \texttt{m'}
with \texttt{m} $\le$ \texttt{m'} and \texttt{P a m'} for the returned value \texttt{a}. The
accompanying combinators are the monad laws for this predicate: \texttt{pure}, \texttt{bind},
\texttt{mono} for weakening the postcondition, \texttt{stateWF} for assuming well-formedness
while proving, \texttt{bind\_le} for carrying an earlier bound, and \texttt{andM} for the
short-circuiting \texttt{<\&\&>}.
\end{definition}

\begin{lemma}[Correctness of the union-find primitives]
\label{fam:eqvmgr-primitives-wf}
\lean{Lean4Lean.EquivManager.find.WF, Lean4Lean.EquivManager.merge.WF, Lean4Lean.EquivManager.toNode.WF}
\leanok
\srcloc{Lean4Lean/Verify/EquivManager.lean}{217}
\uses{def:eqvmgr-mwf, def:kernel-equiv-manager}
\texttt{find i} returns a representative equivalent to \texttt{i} and preserves well-formedness
and the order; \texttt{merge m i j} preserves well-formedness provided the two node ids are the
ids of two expressions already known to be \texttt{IsDefEqE}; and \texttt{toNode e} returns the
node id assigned to \texttt{e}, allocating a fresh one if needed, again preserving
well-formedness.
\end{lemma}
\begin{proof}
\uses{thm:eqvmgr-isdefeqe-relevant, fam:eqvmgr-relevanteq-lemmas, family:expr-axioms}
\leanok
Each unfolds the definition, splits on the bounds check, and rebuilds the two fields of
\texttt{EquivManager.WF} using \texttt{Batteries.UnionFind.equiv\_find}/\texttt{equiv\_union} and
\texttt{Std.HashMap.getElem?\_insert}. \texttt{toNode.WF} is the delicate one: the fresh id is
larger than every existing id, so a claimed equivalence involving it forces the two expressions
to be \texttt{BEq}-equal, hence \texttt{RelevantEq}, hence \texttt{IsDefEqE}.
\end{proof}

\begin{theorem}[Soundness of the cached structural equality test]
\label{thm:eqvmgr-isequiv-wf}
\lean{Lean4Lean.EquivManager.isEquiv.WF}
\leanok
\srcloc{Lean4Lean/Verify/EquivManager.lean}{262}
\uses{def:kernel-equiv-manager, def:eqvmgr-mwf}
Running \texttt{EquivManager.isEquiv useHash e1 e2} from a well-formed manager yields a
well-formed manager, a larger one in the $\le$ order, and, when the boolean result is
\texttt{true}, a derivation of \texttt{IsDefEqE env Us Δ e1 e2}. Only soundness is proved:
nothing says the cache ever answers \texttt{true}, which is the right asymmetry for a kernel.
\end{theorem}
\begin{proof}
\uses{fam:eqvmgr-primitives-wf, axiom:kernel-ptr-eq, family:expr-axioms}
\leanok
Structural recursion mirroring \texttt{isEquiv}: the pointer-equality fast path uses the
\texttt{ptrEqExpr\_eq} axiom (\ref{axiom:kernel-ptr-eq}); the hash mismatch and bound-variable
cases are immediate; otherwise both expressions are interned by \texttt{toNode.WF} and their
representatives found by \texttt{find.WF}, an already-equal pair is discharged by the manager
invariant, and each congruence case recurses and then justifies the \texttt{merge} by
\texttt{merge.WF}. The \texttt{useHash} parameter is threaded but plays no role in soundness.
\end{proof}

\begin{theorem}[Recording a discovered equality preserves the invariant]
\label{thm:eqvmgr-addequiv-wf}
\lean{Lean4Lean.TypeChecker.Inner.addEquiv.WF}
\leanok
\srcloc{Lean4Lean/Verify/EquivManager.lean}{313}
\uses{def:kernel-equiv-manager, ind:trexprs, def:trexpr}
In the type checker's \texttt{RecM}, updating the state by \texttt{EquivManager.addEquiv e1 e2}
preserves the whole \texttt{VState} invariant, provided \texttt{e1} and \texttt{e2} translate to
the same term, one strictly and the other up to definitional equality.
\end{theorem}
\begin{proof}
\uses{fam:eqvmgr-primitives-wf, thm:vtc-methods-iface}
\leanok
Unfold \texttt{addEquiv} into two \texttt{toNode} calls and a \texttt{merge}, then apply
\texttt{toNode.WF} twice and \texttt{merge.WF} once, with the \texttt{IsDefEqE.defeq} witness
obtained by weakening the two translations to the free-variable context.
\end{proof}

\begin{theorem}[Soundness of the memoising \texttt{isDefEq} wrapper]
\label{thm:eqvmgr-isdefeq-wf}
\lean{Lean4Lean.TypeChecker.Inner.isDefEq.WF}
\leanok
\srcloc{Lean4Lean/Verify/EquivManager.lean}{325}
\uses{ind:trexprs, def:hastype-istype}
If \texttt{e1} and \texttt{e2} translate to \texttt{e1'} and \texttt{e2'}, then running the type
checker's \texttt{isDefEq e1 e2} is sound: a \texttt{true} result yields definitional equality of
\texttt{e1'} and \texttt{e2'}, and the resulting state still satisfies the invariant.
\end{theorem}
\begin{proof}
\uses{thm:eqvmgr-addequiv-wf, thm:vtc-methods-iface}
\leanok
Chain the soundness of \texttt{isDefEqCore} with \texttt{addEquiv.WF} for the branch that caches a
positive answer. This is the entry point consumed all over
\texttt{Lean4Lean/Verify/TypeChecker/} and at \srcloc{Lean4Lean/Verify/Environment.lean}{178}.
\end{proof}

\section{Review notes}

Every issue recorded for this group, with severity. No fixes are proposed.

\begin{itemize}
\item \textbf{Medium. Stale doc/code mismatch in the level tests.}
\srcloc{Lean4Lean/Tests/Level.lean}{8} states that what is \emph{not} proved, and so is what the
test file checks, is canonicity of \texttt{normalize'} and completeness of
\texttt{isEquiv'}/\texttt{geq'}. All three are in fact proved:
\srcloc{Lean4Lean/Verify/Level.lean}{3849}, \srcloc{Lean4Lean/Verify/Level.lean}{3861} and
\srcloc{Lean4Lean/Verify/Level.lean}{3870}. The test file's stated purpose no longer matches
reality.
\item \textbf{Medium. Dead code that drags in an extra axiom.} \texttt{mkData\_depth},
\texttt{mkData\_hasParam} and \texttt{mkData\_hasMVar} (\srcloc{Lean4Lean/Verify/Level.lean}{41},
\srcloc{Lean4Lean/Verify/Level.lean}{58}, \srcloc{Lean4Lean/Verify/Level.lean}{73}) are
referenced nowhere in the project. They are the only consumers of the
\texttt{Lean.Level.mkData\_eq} axiom, and they need
\texttt{set\_option allowUnsafeReducibility true} plus
\texttt{attribute [local reducible] Data} (lines 38--39) to elaborate.
\item \textbf{Medium. Trust boundary under the level comparison.} Every main theorem of
\srcloc{Lean4Lean/Verify/Level.lean}{1} factors through \texttt{Std.TreeMap.all\_eq\_all\_toList}
and \texttt{any\_eq\_any\_toList}, declared as axioms at
\srcloc{Lean4Lean/Verify/Axioms.lean}{10} and \srcloc{Lean4Lean/Verify/Axioms.lean}{14} with an
upstream issue link. \texttt{NormLevel.le}, \texttt{NormLevel.addable} and the \texttt{BEq}
instance on \texttt{NormLevel} are all defined through \texttt{TreeMap.all}/\texttt{any}, so the
soundness of the level comparison the kernel performs is modulo these two axioms.
\item \textbf{Medium. The only \texttt{sorryAx} in the group is load-bearing for the memo cache.}
\texttt{RelevantEq.uniq} (\srcloc{Lean4Lean/Verify/EquivManager.lean}{81}) and the two theorems
depending on it reach \texttt{sorryAx} through \texttt{TrProj.uniq}
(\srcloc{Lean4Lean/Verify/Typing/Lemmas.lean}{992}, \texttt{:= sorry}), because the \texttt{proj}
case of \texttt{RelevantEq} compares two projections by index alone and ignores the structure
name. \texttt{isDefEq.WF} (\srcloc{Lean4Lean/Verify/EquivManager.lean}{325}) is the entry point
used throughout \texttt{Verify/TypeChecker/}, so this hole is on the main path of the verified
type checker, not in a corner.
\item \textbf{Low. Stale doc comment in \texttt{LevelStd}.}
\srcloc{Lean4Lean/Verify/LevelStd.lean}{155} still says the two \texttt{qsort} facts are
``currently unproved because \texttt{Array.qsort} has no specification in the standard library''
and names \texttt{qsort\_perm\_toList} and \texttt{pairwise\_qsort\_normLt}, neither of which
exists. Both facts are now proved, at \srcloc{Lean4Lean/Verify/LevelStd.lean}{228} and
\srcloc{Lean4Lean/Verify/LevelStd.lean}{234}, on top of
\srcloc{Lean4Lean/Verify/QSort.lean}{1}.
\item \textbf{Low. A large verified component that nothing consumes.}
\texttt{normalize'\_eval} (\srcloc{Lean4Lean/Verify/Level.lean}{3822}),
\texttt{normalize'\_complete} (\srcloc{Lean4Lean/Verify/Level.lean}{3849}),
\texttt{isEquiv'\_complete} (\srcloc{Lean4Lean/Verify/Level.lean}{3861}),
\texttt{geq'\_wf} (\srcloc{Lean4Lean/Verify/Level.lean}{3828}) and
\texttt{geq'\_complete} (\srcloc{Lean4Lean/Verify/Level.lean}{3870}) are used nowhere. Only
\texttt{isEquivList\_wf} reaches the verified type checker. In particular the whole
reconstruction layer, roughly lines 1290--2250, supports \texttt{normalize'}, which the kernel
never calls: \texttt{geq'} is called at \srcloc{Lean4Lean/Inductive/Add.lean}{226}, whose code
path has no verification yet, and \texttt{normalize'} appears only in
\srcloc{Lean4Lean/Tests/Level.lean}{1}.
\item \textbf{Low. Leftover dead comment block.}
\srcloc{Lean4Lean/Verify/Level.lean}{16} to line 27 contains a commented-out
\texttt{VLevel.toLevel} definition and a \texttt{toLevel\_inj ... := sorry}. Harmless, but it is
the only occurrence of the token \texttt{sorry} in the group and will show up in any audit.
\item \textbf{Low. No module docstring on a 3880-line file.} The internal section docstrings in
\srcloc{Lean4Lean/Verify/Level.lean}{1} are excellent, but a reader reaches line 166 before
learning that the file is about Géran's canonical form; the paper is cited only in
\srcloc{Lean4Lean/Level.lean}{24}.
\item \textbf{Low. Dead code in the equivalence manager.} \texttt{RelevantEq.symm}
(\srcloc{Lean4Lean/Verify/EquivManager.lean}{41}), \texttt{RelevantEq.trans}
(\srcloc{Lean4Lean/Verify/EquivManager.lean}{58}) and \texttt{M.WF.bind\_le}
(\srcloc{Lean4Lean/Verify/EquivManager.lean}{205}) are never used, here or elsewhere; the
similarly named \texttt{RecM.WF.bind\_le} at
\srcloc{Lean4Lean/Verify/TypeChecker/Basic.lean}{343} is a different declaration.
\item \textbf{Low. Overstatement in the \texttt{NormLt} header.}
\srcloc{Lean4Lean/Verify/NormLt.lean}{6} says the file ``shows it is a strict weak order''. What
is exported is asymmetry (\srcloc{Lean4Lean/Verify/NormLt.lean}{350}) and transitivity of the
negation (\srcloc{Lean4Lean/Verify/NormLt.lean}{354}), exactly what
\texttt{Array.qsort\_sorted} takes, not a packaged strict-weak-order statement.
\item \textbf{Low. Vendored code carrying upstream TODOs and cross-namespace side effects.}
\srcloc{Lean4Lean/Verify/QSort.lean}{34} says ``These attributes still need to be moved to the
standard library''; lines 36--51 are commented-out attribute experiments; line 52 says ``These
are just the patterns resulting from \texttt{grind}, but the behaviour should be explained!''.
The file also declares four lemmas at \texttt{\_root\_} into the \texttt{List}, \texttt{Array} and
\texttt{Vector} namespaces (lines 58--103) and registers them, plus several core lemmas, as
global \texttt{grind} lemmas, which every importer inherits.
\item \textbf{Low. Documented gap in the \texttt{qsort} specification.}
\srcloc{Lean4Lean/Verify/QSort.lean}{27} notes that there is no public theorem giving
\texttt{(qsort as lt lo hi)[i] = as[i]} for \texttt{i} outside \texttt{[lo, hi]}. The private
lemmas exist (\srcloc{Lean4Lean/Verify/QSort.lean}{150} and
\srcloc{Lean4Lean/Verify/QSort.lean}{167}) but are not exported; the current consumer needs only
permutation and sortedness, so this is a latent limitation rather than a defect.
\item \textbf{Low. Invariant sprawl in the end-to-end theorems.} Each of them threads five to
seven separate side conditions about \texttt{normalize u} by hand
(\srcloc{Lean4Lean/Verify/Level.lean}{1222} and the invariants at lines 1807, 2413, 2443, 2524
and 3600); \texttt{geq'\_complete} (\srcloc{Lean4Lean/Verify/Level.lean}{3870}) takes six of them
in a single \texttt{refine}. Bundling them into one predicate would make these statements easier
to read and to reuse.
\item \textbf{Low. \texttt{Std.TreeMap} as the representation costs a whole layer.} Because two
maps with the same entries need not be the same tree, every consumer has to be shown to factor
through \texttt{toList} (\ref{fam:lvl-congr}, \srcloc{Lean4Lean/Verify/Level.lean}{3177}). The
design consequence is honestly documented, but a sorted-association-list representation would
have avoided roughly 150 lines of pure bookkeeping and the awkward \texttt{BEq NormLevel} at
\srcloc{Lean4Lean/Verify/Level.lean}{3741}.
\item \textbf{Low. Automation opacity in the two smaller files.}
\texttt{baseCmp\_swap} (\srcloc{Lean4Lean/Verify/NormLt.lean}{97}) closes its last goals with
bare \texttt{grind}, and \texttt{normLtAux\_eq} (\srcloc{Lean4Lean/Verify/NormLt.lean}{264}) ends
several of its fourteen cases in \texttt{simp\_all} or \texttt{grind};
\srcloc{Lean4Lean/Verify/QSort.lean}{1} is \texttt{grind}-first throughout. A human reviewer
cannot check those arguments by reading, and they are brittle against changes in \texttt{grind}.
\item \textbf{Low. Readability of \texttt{baseCmp} and of \texttt{getUndefParam\_none}.}
\texttt{baseCmp} duplicates its \texttt{max} and \texttt{imax} arms verbatim
(\srcloc{Lean4Lean/Verify/NormLt.lean}{47}); the duplication is forced by matching on two
constructors, but it is a readability smell. \texttt{getUndefParam\_none}
(\srcloc{Lean4Lean/Verify/Level.lean}{101}) is written in a dense
\texttt{have ... := ?\_} style with several \texttt{simp}-then-\texttt{split} chains that are
hard to follow.
\item \textbf{Low. Layering oddity in the equivalence manager file.} The last two theorems,
\texttt{addEquiv.WF} (\srcloc{Lean4Lean/Verify/EquivManager.lean}{313}) and \texttt{isDefEq.WF}
(\srcloc{Lean4Lean/Verify/EquivManager.lean}{325}), live in the namespace
\texttt{Lean4Lean.TypeChecker.Inner} rather than \texttt{EquivManager}, and
\texttt{isDefEq.WF} would more naturally sit in
\srcloc{Lean4Lean/Verify/TypeChecker/IsDefEq.lean}{1}. The placement is defensible, since both
depend on the manager invariant, but it is surprising.
\end{itemize}
